\documentclass[11pt, a4paper, oneside]{book}

% ── Engine: XeLaTeX ──────────────────────────────────────────────────────────
\usepackage{fontspec}
\setmainfont{Latin Modern Roman}
\setsansfont{Latin Modern Sans}
% Pin the full system DejaVu Sans Mono by path: the TeX-distribution copy
% is trimmed and silently drops some math glyphs used in the Lean listings
% (e.g. 𝟎, ⬝, ⊕, ∥). The system TTFs carry the full glyph set.
\setmonofont{DejaVuSansMono.ttf}[
  Path=/usr/share/fonts/truetype/dejavu/,
  BoldFont=DejaVuSansMono-Bold.ttf,
  Scale=0.85]

% ── Core packages ────────────────────────────────────────────────────────────
\usepackage[margin=2.5cm]{geometry}
\usepackage{amsmath, amssymb, amsthm, mathtools}
\usepackage{tikz}
\usetikzlibrary{positioning, arrows.meta, calc, decorations.pathmorphing,
  automata, shapes.geometric}
\usepackage{minted}
\setminted{
  fontsize=\small,
  breaklines=true,
  bgcolor=codebg,
  frame=none,
  framesep=2mm,
  baselinestretch=1.15,
}
\usepackage{xcolor}
\usepackage{hyperref}
\hypersetup{colorlinks=true, linkcolor=deepblue, urlcolor=midblue,
  citecolor=accentgreen}
\usepackage{microtype}
\usepackage{graphicx}
\usepackage{enumitem}
\usepackage{booktabs}
\usepackage{tcolorbox}
\tcbuselibrary{theorems, skins, breakable}
\usepackage{fancyhdr}
\usepackage{titlesec}
\usepackage{epigraph}

\usepackage{etoolbox}

% ── Colour palette ───────────────────────────────────────────────────────────
\definecolor{deepblue}{HTML}{1a3a5c}
\definecolor{midblue}{HTML}{2e6da4}
\definecolor{lightblue}{HTML}{d6e8f7}
\definecolor{accentgreen}{HTML}{1e7a4e}
\definecolor{lightgreen}{HTML}{d4edda}
\definecolor{accentpurple}{HTML}{6a3fa0}
\definecolor{lightpurple}{HTML}{ede0f7}
\definecolor{accentorange}{HTML}{b85c00}
\definecolor{lightorange}{HTML}{fde8cc}
\definecolor{accentred}{HTML}{a83232}
\definecolor{lightred}{HTML}{f7d6d6}
\definecolor{codebg}{HTML}{f6f8fa}

% ── Theorem-like environments ────────────────────────────────────────────────
\newtcbtheorem[number within=chapter]{defn}{Definition}{
  enhanced, breakable,
  colback=lightblue, colframe=midblue,
  colupper=black, collower=black,
  fonttitle=\bfseries, coltitle=white,
  separator sign={~---},
}{def}

\newtcbtheorem[number within=chapter]{thm}{Theorem}{
  enhanced, breakable,
  colback=lightgreen, colframe=accentgreen,
  colupper=black, collower=black,
  fonttitle=\bfseries, coltitle=white,
  separator sign={~---},
}{thm}

\newtcbtheorem[number within=chapter]{lem}{Lemma}{
  enhanced, breakable,
  colback=lightgreen!60, colframe=accentgreen!70,
  colupper=black, collower=black,
  fonttitle=\bfseries, coltitle=white,
  separator sign={~---},
}{lem}

\newtcbtheorem[number within=chapter]{prop}{Proposition}{
  enhanced, breakable,
  colback=lightgreen!60, colframe=accentgreen!70,
  colupper=black, collower=black,
  fonttitle=\bfseries, coltitle=white,
  separator sign={~---},
}{prop}

\newtcbtheorem[number within=chapter]{cor}{Corollary}{
  enhanced, breakable,
  colback=lightgreen!40, colframe=accentgreen!50,
  colupper=black, collower=black,
  fonttitle=\bfseries, coltitle=white,
  separator sign={~---},
}{cor}

\newtcolorbox{intuition}[1][]{
  enhanced, breakable,
  colback=lightpurple, colframe=accentpurple,
  colupper=black, fonttitle=\bfseries, coltitle=white,
  title={Intuition}, #1
}

\newtcolorbox{historical}[1][]{
  enhanced, breakable,
  colback=lightorange, colframe=accentorange,
  colupper=black, fonttitle=\bfseries, coltitle=white,
  title={Historical Note}, #1
}

\newtcolorbox{warning}[1][]{
  enhanced, breakable,
  colback=lightred, colframe=accentred,
  colupper=black, fonttitle=\bfseries, coltitle=white,
  title={Warning}, #1
}

\newtcolorbox{exercisebox}[1][]{
  enhanced, breakable,
  colback=lightorange!50, colframe=accentorange!70,
  colupper=black, fonttitle=\bfseries, coltitle=white,
  title={Exercise}, #1
}

\newtcolorbox{leanbox}[1][]{
  enhanced, breakable,
  colback=codebg, colframe=midblue!60,
  colupper=black, fonttitle=\bfseries\ttfamily, coltitle=white,
  title={Lean 4}, #1
}

\newtcolorbox{tsbox}[1][]{
  enhanced, breakable,
  colback=codebg, colframe=accentorange!60,
  colupper=black, fonttitle=\bfseries\ttfamily, coltitle=white,
  title={TypeScript}, #1
}

% ── Custom macros ────────────────────────────────────────────────────────────
\newcommand{\nil}{\mathbf{0}}
\newcommand{\Par}{\mathbin{\|}}
\newcommand{\pref}{\mathbin{.}}
\newcommand{\res}[1]{(\nu #1)}
\newcommand{\repl}{\mathbin{!}}
\newcommand{\inp}[2]{#1(#2)}
\newcommand{\outp}[2]{\overline{#1}\langle #2 \rangle}
\newcommand{\tildep}{\widetilde}
\newcommand{\bisim}{\sim}
\newcommand{\reduces}{\longrightarrow}
\newcommand{\Reduces}{\Longrightarrow}
\newcommand{\fn}{\operatorname{fn}}
\newcommand{\bn}{\operatorname{bn}}
\newcommand{\lc}[1]{\texttt{#1}}
\newcommand{\sem}[1]{\text{⟦}#1\text{⟧}}
\newcommand{\bindnasrepma}{\rotatebox[origin=c]{180}{\&}}

% ── Chapter/section styling ──────────────────────────────────────────────────
\titleformat{\chapter}[display]
  {\normalfont\huge\bfseries\color{deepblue}}
  {\chaptertitlename\ \thechapter}{20pt}{\Huge\color{deepblue}}[\color{black}]
\titleformat{\section}
  {\normalfont\Large\bfseries\color{deepblue}}
  {\thesection}{1em}{}[\color{black}]
\titleformat{\subsection}
  {\normalfont\large\bfseries\color{midblue}}
  {\thesubsection}{1em}{}[\color{black}]

% Global: ensure text is black after any tcolorbox
\AfterEndEnvironment{defn}{\color{black}}
\AfterEndEnvironment{thm}{\color{black}}
\AfterEndEnvironment{lem}{\color{black}}
\AfterEndEnvironment{prop}{\color{black}}
\AfterEndEnvironment{cor}{\color{black}}
\AfterEndEnvironment{intuition}{\color{black}}
\AfterEndEnvironment{historical}{\color{black}}
\AfterEndEnvironment{warning}{\color{black}}
\AfterEndEnvironment{exercisebox}{\color{black}}
\AfterEndEnvironment{leanbox}{\color{black}}
\AfterEndEnvironment{tsbox}{\color{black}}
\AfterEndEnvironment{minted}{\color{black}}

% ── Headers/footers ──────────────────────────────────────────────────────────
\pagestyle{fancy}
\fancyhf{}
\fancyhead[L]{\small\itshape\leftmark}
\fancyhead[R]{\small\thepage}
\renewcommand{\headrulewidth}{0.4pt}

% ── Counter for exercises ────────────────────────────────────────────────────
\newcounter{exercise}[chapter]
\renewcommand{\theexercise}{\thechapter.\arabic{exercise}}

% ══════════════════════════════════════════════════════════════════════════════
\begin{document}

% ── Title page ───────────────────────────────────────────────────────────────
\begin{titlepage}
\centering
\vspace*{3cm}
{\Huge\bfseries\color{deepblue} Process Calculi\\[0.3cm]
\& Verified Concurrency\par}
\vspace{1cm}
{\Large\color{midblue} From Dijkstra's Semaphores to Session Types\\[0.2cm]
in Lean\,4\par}
\vspace{2cm}
{\large A self-contained journey through the theory of\\
concurrent computation, from labeled transition systems\\
to verified deadlock-free protocols.\par}
\vspace{1cm}
{\normalsize With connections to TypeScript/Node.js\\
throughout.\par}
\vspace{3cm}
{\large\itshape Draft --- \today\par}
\vfill
\end{titlepage}

% ── Front matter ─────────────────────────────────────────────────────────────
\frontmatter
\tableofcontents

\chapter*{Preface}
\addcontentsline{toc}{chapter}{Preface}

Concurrent programming is everywhere and almost nobody understands it.

Every web server, every mobile app, every database, every operating system
relies on concurrent execution. Node.js processes thousands of simultaneous
connections through its event loop. Worker threads enable true parallelism in
JavaScript. Go's goroutines and Erlang's actors build entire distributed
systems around message passing. And yet: race conditions, deadlocks, and
data corruption remain among the most common and hardest-to-debug classes of
software defects. The Therac-25 radiation therapy machine killed patients
because of a race condition between concurrent tasks updating shared
state. These are not hypothetical risks.

(A word of caution about famous disasters: not every high-profile software
failure is a concurrency bug. Ariane~5 Flight~501 (1996), for instance, is
often mis-cited as one. It was not: reused Ariane~4 alignment code converted
a 64-bit floating-point horizontal-velocity value into a 16-bit signed
integer, the value overflowed, and the resulting unhandled Operand Error
shut down the guidance computers. That is a data-conversion/overflow defect,
not a race condition. We mention it precisely because misattributing bugs is
itself a hazard---the first step in eliminating a class of defect is naming
it correctly.)

The root cause is that concurrent programs are \emph{fundamentally different}
from sequential ones. When two processes run in parallel, their actions can
interleave in exponentially many ways. Testing can explore only a tiny
fraction of these interleavings. We need \emph{mathematical proof}.

This book develops the theory of concurrent computation from first principles,
formalizes it in Lean\,4, and connects it to the TypeScript/Node.js ecosystem
you use at work. By the end, you will understand:
\begin{itemize}[nosep]
\item What a \emph{process calculus} is and why it's the ``lambda calculus
      of concurrency.''
\item What \emph{bisimulation} means and why it's the right notion of
      equivalence for concurrent systems.
\item How to \emph{prove} that a concurrent protocol is deadlock-free.
\item What \emph{session types} are and how they can guarantee communication
      safety at compile time.
\end{itemize}

\paragraph{Historical arc.}
The intellectual lineage runs: Dijkstra (1965) identified the problem.
Hoare (1978, 1985) and Milner (1980, 1989, 1992) built the theory.
Honda (1993) connected concurrency to type theory. Wadler (2012) connected
it to linear logic. We follow this arc chronologically.

\paragraph{Prerequisites.}
Comfort with Lean\,4 (inductive types, pattern matching, tactics) and
familiarity with TypeScript/Node.js. No prior concurrency theory.

\paragraph{Lean\,4 setup.}
All code is vanilla Lean\,4, no Mathlib---this book is genuinely
zero-dependency, and we build every structure from scratch. The listings are
checked against toolchain \lc{leanprover/lean4:v4.31.0}; a \lc{lean-toolchain}
file containing that single line is all the project needs. Every printed
listing compiles as shown (the one exception, an intentional \lc{sorry}, is an
exercise and is flagged where it appears). Do not worry if the first listing
is your first encounter with Lean~4: Appendix~\ref{app:primer}
(\emph{A Lean 4 Primer}) explains every keyword and tactic this book uses,
with small worked examples and---where it illuminates---the lower-level term
each tactic expands to. Read it alongside the early chapters, or dip in
whenever a new construct appears.

\mainmatter
\color{black}

% ══════════════════════════════════════════════════════════════════════════════
\part{Foundations of Communicating Processes}
% ══════════════════════════════════════════════════════════════════════════════

% ══════════════════════════════════════════════════════════════════════════════
\chapter{Why Concurrency Is Hard}
% ══════════════════════════════════════════════════════════════════════════════

\epigraph{``Testing can be used to show the presence of bugs, but never to show
their absence.''}{\textit{Edsger W.\ Dijkstra}}

\section{The Problem}

\begin{historical}[title={Dijkstra and the Birth of Concurrency Theory}]
In 1965, Edsger Dijkstra published ``Cooperating Sequential Processes,'' a
technical report that identified the fundamental problem of concurrent
programming: when multiple processes share resources, things go wrong in ways
that are almost impossible to detect by testing. A sequential program always
does the same thing given the same input. A concurrent program can do
\emph{different} things on different runs, depending on the precise timing of
events---the \emph{interleaving} of actions.

Dijkstra introduced the \textbf{mutual exclusion problem}: given $n$ processes
that need occasional access to a shared resource, design a protocol ensuring
that at most one process uses the resource at any time. His solution used
\textbf{semaphores}---integer variables with atomic increment and decrement
operations---the first general-purpose synchronization primitive. His 1968
paper ``Go To Statement Considered Harmful'' made him famous for structured
programming, but his concurrency work is arguably his deepest contribution.

Dijkstra received the Turing Award in 1972. His acceptance lecture contained
the memorable line: ``The question of whether machines can think is about as
relevant as the question of whether submarines can swim.''
\end{historical}

\section{Concurrency in Node.js}

As a TypeScript/Node.js programmer, you deal with concurrency every day---even
if you don't always call it that.

\subsection{The Event Loop}

Node.js is single-threaded but \emph{asynchronous}. The event loop processes
callbacks one at a time, but I/O operations run in the background:

\begin{tsbox}
\begin{minted}{typescript}
// These run "concurrently" --- the event loop interleaves them
const fetchUser = async (id: string) => {
  const resp = await fetch(`/api/users/${id}`);  // yields to event loop
  return resp.json();
};

// Two concurrent fetches
const [alice, bob] = await Promise.all([
  fetchUser("alice"),
  fetchUser("bob"),
]);
\end{minted}
\end{tsbox}

This is \emph{cooperative concurrency}: each task voluntarily yields at
\lc{await} points. No two callbacks run simultaneously, so you avoid many
classical race conditions. But not all of them---the interleaving of
callbacks can still surprise you.

\subsection{Worker Threads}

For CPU-intensive work, Node.js provides \lc{Worker} threads---true
parallelism:

\begin{tsbox}
\begin{minted}{typescript}
import { Worker, parentPort, MessageChannel } from 'worker_threads';

// Main thread
const worker = new Worker('./worker.js');
worker.postMessage({ task: 'compute', data: payload });
worker.on('message', (result) => {
  console.log('Got result:', result);
});

// Worker thread (worker.js)
parentPort?.on('message', (msg) => {
  const result = heavyComputation(msg.data);
  parentPort?.postMessage(result);
});
\end{minted}
\end{tsbox}

The moment you use workers, you enter Dijkstra's world: two threads running
\emph{truly simultaneously}, communicating via message passing. The
\lc{postMessage}/\lc{on('message')} pattern is exactly what process calculi
formalize.

\section{What Can Go Wrong}

\subsection{Race Conditions}

\begin{tsbox}
\begin{minted}{typescript}
// Shared state between event loop callbacks
let balance = 100;

async function withdraw(amount: number) {
  if (balance >= amount) {        // Check
    await delay(10);              // Simulate async work
    balance -= amount;            // Act --- but balance may have changed!
  }
}

// Two concurrent withdrawals: both see balance = 100
await Promise.all([withdraw(80), withdraw(80)]);
// balance is now -60, not 20!
\end{minted}
\end{tsbox}

\subsection{Deadlock}

\begin{tsbox}
\begin{minted}{typescript}
// Worker A waits for B's response before sending its own
// Worker B waits for A's response before sending its own
// Neither can proceed: DEADLOCK

workerA.postMessage('waiting for B');  // A blocks until B responds
workerB.postMessage('waiting for A');  // B blocks until A responds
\end{minted}
\end{tsbox}

\begin{historical}[title={The Therac-25}]
The Therac-25 (1985--1987) was a radiation therapy machine whose software
had a race condition between the operator interface and the beam control
system. Under specific timing---an operator typing faster than the system
expected---the machine delivered radiation doses hundreds of times higher
than intended. Six patients were seriously injured; three died. The bug was
virtually impossible to reproduce in testing because it depended on exact
keystroke timing. This is the canonical example of why concurrent systems
need mathematical reasoning, not just testing.
\end{historical}

\section{The Way Forward}

The rest of this book develops mathematical tools for reasoning about
concurrent systems with \emph{certainty}:
\begin{itemize}[nosep]
\item \textbf{Labeled transition systems} (Chapter~2): model concurrent
  systems as state machines.
\item \textbf{CCS} (Chapter~3): a ``programming language'' for concurrent
  processes.
\item \textbf{Bisimulation} (Chapter~4): when are two concurrent systems
  equivalent?
\item \textbf{Session types} (Chapter~9): type systems that guarantee
  deadlock freedom.
\end{itemize}

Every concept will be formalized in Lean\,4 and connected back to
TypeScript.

% ══════════════════════════════════════════════════════════════════════════════
\chapter{Labeled Transition Systems}
% ══════════════════════════════════════════════════════════════════════════════

\epigraph{``The nice thing about standards is that you have so many to choose
from.''}{\textit{Andrew S.\ Tanenbaum}}

\section{Processes as State Machines}

\begin{historical}[title={Keller and the LTS}]
Robert Keller (1976) introduced labeled transition systems as a mathematical
semantics for concurrent programs. The idea: a process is a state machine
whose transitions are labeled with \emph{actions}---things the process does
that the outside world can observe---or with a special silent action $\tau$
(something internal the observer cannot see). This framework became the
standard foundation for all subsequent process calculi.
\end{historical}

\begin{defn}{Labeled Transition System}{lts}
A \textbf{labeled transition system (LTS)} is a triple $(S, A, \to)$ where:
\begin{itemize}[nosep]
\item $S$ is a set of \textbf{states}.
\item $A$ is a set of \textbf{actions} (labels), including the
  \textbf{silent action} $\tau$.
\item ${\to} \subseteq S \times A \times S$ is the \textbf{transition
  relation}. We write $s \xrightarrow{a} s'$ for $(s, a, s') \in {\to}$.
\end{itemize}
\end{defn}

\begin{leanbox}
\begin{minted}{lean}
-- Actions: either observable or silent (tau)
inductive Action (Label : Type) where
  | obs : Label → Action Label    -- observable action
  | tau : Action Label             -- internal/silent action
  deriving DecidableEq, Repr

-- A labeled transition system
structure LTS (State Label : Type) where
  step : State → Action Label → State → Prop
\end{minted}
\end{leanbox}

\section{Traces}

A \textbf{trace} of a process is a sequence of \emph{observable} actions it
can perform (ignoring $\tau$ steps). Formally, the trace set of state $s$ is:
\[
  \mathrm{traces}(s) = \{a_1 a_2 \cdots a_n :
    \exists s_0, \ldots, s_n,\;
    s = s_0 \Reduces^{*} \xrightarrow{a_1} \Reduces^{*}
    \xrightarrow{a_2} \cdots \xrightarrow{a_n} s_n \}
\]
where $\Reduces^{*}$ denotes zero or more $\tau$ steps.

\section{Example: A Vending Machine}

\begin{center}
\begin{tikzpicture}[
    ->, >=stealth, node distance=2.5cm,
    state/.style={circle, draw, fill=lightblue!50, minimum size=1cm,
      font=\small},
  ]
  \node[state] (idle) {idle};
  \node[state, right of=idle] (paid) {paid};
  \node[state, above right=1cm and 2cm of paid] (coffee) {$\checkmark$};
  \node[state, below right=1cm and 2cm of paid] (tea) {$\checkmark$};

  \path (idle)  edge node[above] {\small coin} (paid)
        (paid)  edge node[above left] {\small coffee} (coffee)
        (paid)  edge node[below left] {\small tea} (tea);
\end{tikzpicture}
\end{center}

This machine accepts a coin, then dispenses either coffee or tea. Its trace
set is $\{\texttt{coin}\cdot\texttt{coffee},\; \texttt{coin}\cdot\texttt{tea}\}$.

\begin{leanbox}
\begin{minted}{lean}
-- Vending machine states
inductive VMState where
  | idle | paid | done
  deriving DecidableEq, Repr

-- Vending machine actions
inductive VMLabel where
  | coin | coffee | tea
  deriving DecidableEq, Repr

-- Transition relation
inductive VMStep : VMState → Action VMLabel → VMState → Prop where
  | insert : VMStep .idle (.obs .coin) .paid
  | getCoffee : VMStep .paid (.obs .coffee) .done
  | getTea : VMStep .paid (.obs .tea) .done
\end{minted}
\end{leanbox}

\section{Determinism and Nondeterminism}

The vending machine is \textbf{nondeterministic}: from state \lc{paid},
there are two possible transitions (coffee or tea). The ``choice'' is made
by the environment (the customer), not the machine.

\begin{defn}{Deterministic LTS}{det-lts}
An LTS is \textbf{deterministic} if for every state $s$ and action $a$,
there is at most one state $s'$ with $s \xrightarrow{a} s'$.
\end{defn}

Most interesting concurrent systems are nondeterministic: the interleaving
of parallel processes creates multiple possible execution paths.

\begin{exercisebox}[title={Exercise 2.1}]
\stepcounter{exercise}
Model a traffic light (green $\to$ yellow $\to$ red $\to$ green) as an LTS
in Lean. Prove it is deterministic.
\end{exercisebox}

\begin{exercisebox}[title={Exercise 2.2}]
\stepcounter{exercise}
Model a retry-with-backoff mechanism: a process that sends a request, either
receives a response (success) or times out (failure), and on failure retries
up to 3 times. Define the LTS and compute its trace set.
\end{exercisebox}

% ══════════════════════════════════════════════════════════════════════════════
\chapter{CCS --- The Calculus of Communicating Systems}
% ══════════════════════════════════════════════════════════════════════════════

\epigraph{``There is no concurrency theory without Robin Milner.''}%
{\textit{Luca Cardelli}}

\section{Robin Milner and the Quest for a Theory}

\begin{historical}[title={Robin Milner (1934--2010)}]
Robin Milner is one of the most important computer scientists in history. He
invented three things, each of which alone would be a career-defining
achievement:

\textbf{ML} (1973): the programming language family that includes OCaml and
F\#, and whose type inference algorithm (Algorithm W) inspired Haskell and
Lean's type systems. Milner proved that ``well-typed programs do not go
wrong''---the first type soundness theorem.

\textbf{CCS} (1980): the Calculus of Communicating Systems, a minimal formal
language for concurrent processes. Just as Church's lambda calculus captures
sequential computation with three constructs (variable, abstraction,
application), CCS captures concurrent computation with a small set of
operators (inaction, prefix, choice, parallel, restriction).

\textbf{The $\pi$-calculus} (1992, with Parrow and Walker): an extension of
CCS where channel names can be communicated, modeling systems whose
communication topology changes at runtime.

Milner received the Turing Award in 1991 ``for three distinct and complete
achievements''---an unprecedented citation. His 1989 book
\textit{Communication and Concurrency} remains the definitive introduction
to the theory.
\end{historical}

\section{CCS Syntax}

CCS processes are built from six constructors:

\begin{defn}{CCS Syntax}{ccs-syntax}
Given a set of \textbf{names} (channels) $\mathcal{N}$, the set of CCS
processes is defined by:
\[
  P, Q \;::=\;
    \nil
    \;\mid\; \alpha \pref P
    \;\mid\; P + Q
    \;\mid\; P \Par Q
    \;\mid\; \res{a} P
    \;\mid\; \repl P
\]
where $\alpha$ is an \textbf{action}: either a name $a$ (input), a co-name
$\bar{a}$ (output), or $\tau$ (silent).
\end{defn}

\paragraph{What each constructor means:}
\begin{itemize}[nosep]
\item $\nil$ --- \textbf{inaction}. The stopped process; does nothing.
\item $\alpha \pref P$ --- \textbf{prefix}. Perform action $\alpha$, then
  continue as $P$.
\item $P + Q$ --- \textbf{choice}. Behave as $P$ or $Q$; the environment
  decides which.
\item $P \Par Q$ --- \textbf{parallel composition}. $P$ and $Q$ run
  concurrently.
\item $\res{a} P$ --- \textbf{restriction}. Channel $a$ is private to $P$;
  the outside world cannot use it.
\item $\repl P$ --- \textbf{replication}. Unbounded copies of $P$;
  models a server that can handle any number of requests.
\end{itemize}

\begin{leanbox}
\begin{minted}{lean}
-- Channel names
abbrev Name := String

-- CCS actions
inductive CCSAct where
  | inp : Name → CCSAct          -- input on channel a
  | out : Name → CCSAct          -- output on channel ā
  | tau : CCSAct                  -- silent action
  deriving DecidableEq, Repr

-- CCS is the complement of an action
def CCSAct.complement : CCSAct → Option CCSAct
  | .inp a => some (.out a)
  | .out a => some (.inp a)
  | .tau   => none

-- CCS processes
inductive CCS where
  | nil  : CCS                           -- 0
  | pre  : CCSAct → CCS → CCS           -- α.P
  | plus : CCS → CCS → CCS              -- P + Q
  | par  : CCS → CCS → CCS              -- P | Q
  | res  : Name → CCS → CCS             -- (νa)P
  | repl : CCS → CCS                    -- !P
  deriving Repr

-- Notation helpers. `scoped` attributes are only legal inside a
-- namespace, so we wrap them (open CCSNotation to bring them in scope).
namespace CCSNotation
scoped notation "∅" => CCS.nil
scoped infixr:60 " ⊳ " => CCS.pre        -- α.P
scoped infixl:50 " ⊕ " => CCS.plus       -- P + Q
scoped infixl:40 " ‖ " => CCS.par        -- P | Q
end CCSNotation
\end{minted}
\end{leanbox}

\section{Operational Semantics}

The \emph{meaning} of a CCS process is given by its operational semantics:
an inductive relation defining which transitions are possible.

\begin{defn}{CCS Operational Semantics}{ccs-sos}
The transition relation $P \xrightarrow{\alpha} P'$ is the smallest relation
satisfying these rules:
\begin{enumerate}[nosep, label=(\textsc{\alph*})]
\item \textsc{Prefix:} $\alpha \pref P \xrightarrow{\alpha} P$
\item \textsc{Choice-L:} If $P \xrightarrow{\alpha} P'$ then
  $P + Q \xrightarrow{\alpha} P'$
\item \textsc{Choice-R:} If $Q \xrightarrow{\alpha} Q'$ then
  $P + Q \xrightarrow{\alpha} Q'$
\item \textsc{Par-L:} If $P \xrightarrow{\alpha} P'$ then
  $P \Par Q \xrightarrow{\alpha} P' \Par Q$
\item \textsc{Par-R:} If $Q \xrightarrow{\alpha} Q'$ then
  $P \Par Q \xrightarrow{\alpha} P \Par Q'$
\item \textsc{Comm:} If $P \xrightarrow{a} P'$ and $Q \xrightarrow{\bar{a}} Q'$
  then $P \Par Q \xrightarrow{\tau} P' \Par Q'$ (and symmetrically, with the
  roles of $P$ and $Q$ swapped---the label $a$ ranges over both inputs and
  their complementary outputs, so a send on the left meeting a receive on the
  right also synchronizes). The Lean encoding below makes both directions
  explicit as the constructors \lc{comm} and \lc{commR}.
\item \textsc{Res:} If $P \xrightarrow{\alpha} P'$ and $\alpha \neq a, \bar{a}$
  then $\res{a} P \xrightarrow{\alpha} \res{a} P'$
\item \textsc{Repl:} If $P \Par \repl P \xrightarrow{\alpha} P'$
  then $\repl P \xrightarrow{\alpha} P'$
\end{enumerate}
\end{defn}

\begin{leanbox}
\begin{minted}{lean}
-- CCS transition relation
inductive CCSStep : CCS → CCSAct → CCS → Prop where
  | prefix : CCSStep (.pre α P) α P
  | choiceL : CCSStep P α P' → CCSStep (.plus P Q) α P'
  | choiceR : CCSStep Q α Q' → CCSStep (.plus P Q) α Q'
  | parL : CCSStep P α P' → CCSStep (.par P Q) α (.par P' Q)
  | parR : CCSStep Q α Q' → CCSStep (.par P Q) α (.par P Q')
  | comm : CCSStep P (.inp a) P' → CCSStep Q (.out a) Q' →
           CCSStep (.par P Q) .tau (.par P' Q')
  | commR : CCSStep P (.out a) P' → CCSStep Q (.inp a) Q' →
            CCSStep (.par P Q) .tau (.par P' Q')
  | res : CCSStep P α P' → α ≠ .inp a → α ≠ .out a →
          CCSStep (.res a P) α (.res a P')
\end{minted}
\end{leanbox}

\section{The Communication Rule}

The most important rule is \textsc{Comm}: when two parallel processes perform
complementary actions ($a$ and $\bar{a}$) on the same channel, they
\textbf{synchronize} via a silent $\tau$ step. This models message passing:
one process sends, the other receives, and the communication is invisible to
the outside.

\subsection{Worked Example: Request-Response}

\begin{align*}
  \text{Client} &= \overline{\text{req}} \pref \text{resp} \pref \nil \\
  \text{Server} &= \text{req} \pref \overline{\text{resp}} \pref \nil \\[6pt]
  \text{Client} \Par \text{Server}
  &\xrightarrow{\tau} \text{resp} \pref \nil \;\Par\;
    \overline{\text{resp}} \pref \nil
  &&\text{(\textsc{Comm}: req/}\overline{\text{req}}\text{)} \\
  &\xrightarrow{\tau} \nil \Par \nil
  &&\text{(\textsc{Comm}: resp/}\overline{\text{resp}}\text{)} \\
  &= \nil \Par \nil
  &&\text{(done)}
\end{align*}

Two $\tau$ steps: two internal communications. The outside world sees nothing
---the entire exchange is private.

\subsection{The Node.js Mapping}

\begin{center}
\begin{tabular}{lcl}
  \textbf{CCS} & & \textbf{Node.js} \\
  \midrule
  $P \Par Q$ & $\longleftrightarrow$ & two \lc{Worker} threads \\
  $a$ / $\bar{a}$ & $\longleftrightarrow$ & \lc{port.on('message')} / \lc{port.postMessage()} \\
  $\res{a} P$ & $\longleftrightarrow$ & \lc{new MessageChannel()} local to a module \\
  $P + Q$ & $\longleftrightarrow$ & \lc{Promise.race([p, q])} \\
  $\repl P$ & $\longleftrightarrow$ & server listening in a loop \\
\end{tabular}
\end{center}

\begin{warning}[title={The \lc{Promise.race} analogy is imperfect}]
CCS choice $P + Q$ \emph{discards} the branch not taken: the moment $P$ takes
a step, $Q$ is gone, and none of $Q$'s effects ever happen.
\lc{Promise.race([p, q])} resolves with whichever promise settles first, but
it does \textbf{not} cancel the loser---the losing promise keeps running and
its side effects still occur. So the mapping captures ``first one wins the
result,'' not ``the other is abandoned.'' A closer operational match to $+$ is
an \lc{AbortController} whose \lc{abort()} genuinely tears down the loser.
\end{warning}

\begin{exercisebox}[title={Exercise 3.1}]
\stepcounter{exercise}
Model a pub/sub system in CCS: a publisher sends messages on a channel, and
two subscribers receive them. Define the processes and trace through the
possible executions.
\end{exercisebox}

\begin{exercisebox}[title={Exercise 3.2}]
\stepcounter{exercise}
Prove in Lean that the Client $\|$ Server system above reaches
$\nil \Par \nil$ after exactly two $\tau$ steps, by constructing the
derivation using \lc{CCSStep} constructors. \emph{Hint:} each step is a single
constructor --- \lc{CCSStep.commR CCSStep.prefix CCSStep.prefix} for the first
(Client outputs, Server inputs) and \lc{CCSStep.comm …} for the second. If
tactic and term proofs are new to you, see Appendix~\ref{app:primer}; a fully
worked version of this derivation is in the companion file
\lc{PiCalcVerification.lean}.
\end{exercisebox}

% ══════════════════════════════════════════════════════════════════════════════
\chapter{Bisimulation --- When Are Two Systems ``The Same''?}
% ══════════════════════════════════════════════════════════════════════════════

\epigraph{``Two systems are equivalent if no experiment can distinguish
them.''}{\textit{Robin Milner} (paraphrased)}

\section{The Coffee Machine Problem}

\begin{historical}[title={Park and Milner}]
After defining CCS, Milner faced a subtle question: when should two
processes be considered equivalent? The naive answer---``they can perform
the same sequences of actions'' (trace equivalence)---turns out to be too
coarse. David Park (1981) and Milner independently arrived at the right
notion: \textbf{bisimulation}. Two processes are bisimilar if every step one
can take, the other can match---step by step, preserving the branching
structure, not just the complete traces.
\end{historical}

The canonical counterexample:
\begin{align*}
  P &= \text{coin} \pref (\text{coffee} \pref \nil + \text{tea} \pref \nil) \\
  Q &= \text{coin} \pref \text{coffee} \pref \nil
       + \text{coin} \pref \text{tea} \pref \nil
\end{align*}

Both have the same traces: $\{\text{coin}\cdot\text{coffee},\;
\text{coin}\cdot\text{tea}\}$. But they behave differently! In $P$, after
inserting a coin you can still choose coffee or tea. In $Q$, the choice was
already made (by the nondeterministic resolution of $+$) \emph{before} you
inserted the coin. An observer who inserts a coin and then tries to get
coffee might fail with $Q$ (if the second branch was chosen) but never with
$P$.

\begin{center}
\begin{tikzpicture}[
    ->, >=stealth, node distance=1.8cm,
    state/.style={circle, draw, fill=lightblue!50, minimum size=0.7cm,
      font=\footnotesize},
  ]
  % P
  \node[state] (P0) {$P$};
  \node[state, below of=P0] (P1) {$P_1$};
  \node[state, below left=1.2cm of P1] (Pc) {$\checkmark$};
  \node[state, below right=1.2cm of P1] (Pt) {$\checkmark$};
  \path (P0) edge node[right, font=\footnotesize] {coin} (P1)
        (P1) edge node[left, font=\footnotesize] {coffee} (Pc)
        (P1) edge node[right, font=\footnotesize] {tea} (Pt);
  \node[below=0.3cm of Pc, font=\small\bfseries] (Plbl) {};
  \node[above=0.3cm of P0, font=\small\bfseries\color{deepblue}] {$P$: choice after coin};

  % Q
  \node[state, right=4cm of P0] (Q0) {$Q$};
  \node[state, below left=1.5cm of Q0] (Q1) {};
  \node[state, below right=1.5cm of Q0] (Q2) {};
  \node[state, below of=Q1] (Qc) {$\checkmark$};
  \node[state, below of=Q2] (Qt) {$\checkmark$};
  \path (Q0) edge node[left, font=\footnotesize] {coin} (Q1)
        (Q0) edge node[right, font=\footnotesize] {coin} (Q2)
        (Q1) edge node[left, font=\footnotesize] {coffee} (Qc)
        (Q2) edge node[right, font=\footnotesize] {tea} (Qt);
  \node[above=0.3cm of Q0, font=\small\bfseries\color{deepblue}] {$Q$: choice before coin};
\end{tikzpicture}
\end{center}

\section{Bisimulation}

\begin{defn}{Bisimulation}{bisimulation}
Let $(S, A, \to)$ be an LTS. A binary relation $R \subseteq S \times S$ is a
\textbf{bisimulation} if whenever $s \mathrel{R} t$:
\begin{enumerate}[nosep, label=(\roman*)]
\item If $s \xrightarrow{a} s'$, then there exists $t'$ such that
  $t \xrightarrow{a} t'$ and $s' \mathrel{R} t'$.
\item If $t \xrightarrow{a} t'$, then there exists $s'$ such that
  $s \xrightarrow{a} s'$ and $s' \mathrel{R} t'$.
\end{enumerate}
Two states $s, t$ are \textbf{bisimilar} (written $s \bisim t$) if there
exists a bisimulation $R$ with $s \mathrel{R} t$.
\end{defn}

\begin{leanbox}
\begin{minted}{lean}
-- A relation R is a bisimulation
def IsBisimulation {State Label : Type}
    (step : State → Action Label → State → Prop)
    (R : State → State → Prop) : Prop :=
  ∀ s t, R s t →
    (∀ a s', step s a s' → ∃ t', step t a t' ∧ R s' t') ∧
    (∀ a t', step t a t' → ∃ s', step s a s' ∧ R s' t')

-- Bisimilarity: exists a bisimulation relating them
def Bisimilar {State Label : Type}
    (step : State → Action Label → State → Prop)
    (s t : State) : Prop :=
  ∃ R, IsBisimulation step R ∧ R s t
\end{minted}
\end{leanbox}

\section{Bisimulation as a Greatest Fixed Point}

\begin{intuition}
Bisimilarity has a beautiful lattice-theoretic characterization. Define a
function $F$ on relations: $F(R) = \{(s,t) \mid \text{$R$ satisfies the
bisimulation conditions at $(s,t)$}\}$. Then bisimilarity is the
\textbf{greatest fixed point} of $F$---the largest relation that is a
bisimulation.

This is the Knaster-Tarski theorem from lattice theory! The lattice of
relations on $S$ (ordered by inclusion) is complete, $F$ is monotone, and
its greatest fixed point exists and equals the union of all bisimulations.

If you've studied the lattice/propagator book, this connection is direct:
bisimilarity is computed by the same fixed-point machinery that computes
propagator convergence.
\end{intuition}

\section{Properties of Bisimilarity}

\begin{thm}{Bisimilarity is an Equivalence Relation}{bisim-equiv}
$\bisim$ is reflexive, symmetric, and transitive.
\end{thm}

\begin{proof}
Reflexive: the identity relation is a bisimulation. Symmetric: if $R$ is a
bisimulation, so is $R^{-1}$. Transitive: if $R_1$ and $R_2$ are
bisimulations, then $R_1 \circ R_2$ is a bisimulation (this requires showing
that the composed relation satisfies both conditions).
\end{proof}

\begin{thm}{Bisimilarity is a Congruence}{bisim-congruence}
If $P \bisim Q$, then:
$\alpha \pref P \bisim \alpha \pref Q$, \quad
$P + R \bisim Q + R$, \quad
$P \Par R \bisim Q \Par R$, \quad
$\res{a} P \bisim \res{a} Q$.
\end{thm}

This means bisimilar processes can be substituted for each other in any
context---the analogue of the ``substitution principle'' in programming.

\begin{exercisebox}[title={Exercise 4.1}]
\stepcounter{exercise}
Prove that $P \Par Q \bisim Q \Par P$ (parallel composition is commutative
up to bisimulation). Construct the bisimulation relation explicitly.
\end{exercisebox}

\begin{exercisebox}[title={Exercise 4.2}]
\stepcounter{exercise}
Prove that $P$ and $Q$ (the coffee machines above) are \emph{not} bisimilar.
\emph{Hint:} after the coin step, $P$ reaches a state that can do both
coffee and tea, but $Q$ reaches a state that can do only one.
\end{exercisebox}

% ══════════════════════════════════════════════════════════════════════════════
\chapter{Deadlock, Livelock, and Safety Properties}
% ══════════════════════════════════════════════════════════════════════════════

\epigraph{``Simplicity is a great virtue but it requires hard work to achieve
it and education to appreciate it.''}{\textit{Edsger W.\ Dijkstra}}

\section{Dining Philosophers}

\begin{historical}[title={Dijkstra's Dining Philosophers (1971)}]
Five philosophers sit around a circular table. Between each pair is a single
fork. To eat, a philosopher needs both adjacent forks. Each philosopher
alternates between thinking and eating. If every philosopher simultaneously
picks up their left fork, then all wait forever for the right fork:
\textbf{deadlock}.

Dijkstra introduced this problem to illustrate that deadlock is not a bug in
any individual component---each philosopher follows a perfectly reasonable
protocol. Deadlock is a property of the \emph{interaction}, arising from
circular dependencies.
\end{historical}

\section{Deadlock}

\begin{defn}{Deadlock}{deadlock}
A CCS process $P$ is \textbf{deadlocked} if:
\begin{enumerate}[nosep, label=(\roman*)]
\item $P \neq \nil$ (the process hasn't terminated).
\item There is no $\alpha$ and $P'$ such that $P \xrightarrow{\alpha} P'$
  (no transitions possible).
\end{enumerate}
\end{defn}

\begin{leanbox}
\begin{minted}{lean}
def isDeadlocked (P : CCS) : Prop :=
  P ≠ .nil ∧ ∀ α P', ¬ CCSStep P α P'
\end{minted}
\end{leanbox}

\subsection{Deadlock in CCS}

\begin{align*}
  A &= \overline{b} \pref a \pref \nil &&\text{(sends on $b$, then waits on $a$)} \\
  B &= \overline{a} \pref b \pref \nil &&\text{(sends on $a$, then waits on $b$)} \\[4pt]
  &\res{a}\res{b}(A \Par B)
\end{align*}

Both $A$ and $B$ want to \emph{send} first. Under restriction, neither output
is visible externally, so a step could only happen by an internal
synchronization---and none is possible. $A$'s first (and blocking) action is
an \emph{output} on $b$, but nobody is listening on $b$: $B$'s first action is
an output on $a$, not an input on $b$. Symmetrically, $A$ never offers an
input on $a$ to match $B$'s output. Neither process ever reaches its later
input action, because each is stuck trying to send. Two senders, no matching
receivers. Deadlock.

\section{Livelock}

\begin{defn}{Livelock / Divergence}{livelock}
A process $P$ \textbf{diverges} (livelocks) if it can perform an infinite
sequence of $\tau$ steps without ever performing an observable action.
\end{defn}

In TypeScript terms: a \lc{setImmediate} loop that prevents the event loop
from processing any other callbacks.

\section{Safety and Liveness}

\begin{defn}{Safety Property}{safety}
A \textbf{safety property} asserts: ``nothing bad ever happens.'' Formally,
a predicate $\varphi$ holds in all reachable states.
\end{defn}

\begin{defn}{Liveness Property}{liveness}
A \textbf{liveness property} asserts: ``something good eventually happens.''
Formally, every execution eventually reaches a state satisfying $\varphi$.
\end{defn}

\begin{historical}[title={Hoare and CSP}]
Tony Hoare's \textit{Communicating Sequential Processes} (1978 paper, 1985
book) is the other major process calculus, developed independently from
Milner's CCS. Hoare focused on \emph{traces and refinement}---specifying what
a system should do and verifying that an implementation refines the
specification. CSP directly influenced Go's goroutines and channels; Rob Pike,
one of Go's creators, cites Hoare as the primary inspiration.

CSP and CCS are closely related but have different philosophies: CSP focuses
on traces (good for specification), CCS on bisimulation (good for
equivalence). Both are Turing-complete and can encode each other.
\end{historical}

\section{The Dining Philosophers in CCS}

\begin{leanbox}
\begin{minted}{lean}
-- Philosopher i picks up fork i, then fork (i+1) mod 5, eats, puts both down
def philosopher (i : Fin 5) : CCS :=
  let li := s!"pick{i}"
  let ri := s!"pick{(i.val + 1) % 5}"
  let di := s!"drop{i}"
  let dri := s!"drop{(i.val + 1) % 5}"
  .pre (.inp li) <|     -- pick up left fork
  .pre (.inp ri) <|     -- pick up right fork
  .pre (.out di) <|     -- put down left fork
  .pre (.out dri) <|    -- put down right fork
  .nil

-- The asymmetric fix: philosopher 4 picks up forks in reverse order
def philosopher4_fixed : CCS :=
  .pre (.inp "pick0") <|     -- right fork first!
  .pre (.inp "pick4") <|     -- then left fork
  .pre (.out "drop0") <|
  .pre (.out "drop4") <|
  .nil
\end{minted}
\end{leanbox}

The asymmetric solution breaks the circular dependency: philosopher~4 picks up
fork~0 first (instead of fork~4), which means at least one philosopher can
always make progress.

\begin{exercisebox}[title={Exercise 5.1}]
\stepcounter{exercise}
Model a reader-writer lock in CCS: multiple readers can hold the lock
simultaneously, but a writer needs exclusive access. Prove (informally, or
in Lean) that the protocol is deadlock-free.
\end{exercisebox}

\begin{exercisebox}[title={Exercise 5.2}]
\stepcounter{exercise}
Formalize the deadlock property in Lean for the symmetric dining philosophers
(all pick up left fork first). Show that a deadlocked configuration exists.
\end{exercisebox}

% ══════════════════════════════════════════════════════════════════════════════
\part{Advanced Theory}
% ══════════════════════════════════════════════════════════════════════════════

% ══════════════════════════════════════════════════════════════════════════════
\chapter{The $\pi$-Calculus --- Mobile Channels}
% ══════════════════════════════════════════════════════════════════════════════

\epigraph{``Functions are processes, and processes are functions.''}%
{\textit{Robin Milner}}

\section{Beyond Fixed Topologies}

\begin{historical}[title={Milner, Parrow, and Walker (1992)}]
CCS has a limitation: the set of channel names is fixed at design time. In
real distributed systems, communication topologies change at runtime: a
server hands off a client connection to a worker; a chat participant invites
a new user; a microservice discovers another via a registry.

In 1992, Milner (with Joachim Parrow and David Walker) introduced the
\textbf{$\pi$-calculus}, extending CCS with one radical new feature:
\textbf{channel names can be communicated over channels}. This seems like a
small change---$\outp{x}{y}.P$ means ``send the name $y$ over channel $x$,
then continue as $P$''---but its consequences are profound. It models any
system whose communication graph evolves dynamically.

Milner's tour de force was showing that the $\lambda$-calculus can be
\emph{encoded} in the $\pi$-calculus: function application becomes
communication, closures become restricted channels. This means concurrent
computation is at least as fundamental as sequential computation.
\end{historical}

\section{$\pi$-Calculus Syntax}

\begin{defn}{$\pi$-Calculus}{pi-syntax}
\[
  P, Q \;::=\;
    \nil
    \;\mid\; \outp{x}{y} \pref P
    \;\mid\; \inp{x}{y} \pref P
    \;\mid\; P \Par Q
    \;\mid\; \res{a} P
    \;\mid\; \repl P
    \;\mid\; P + Q
\]
\begin{itemize}[nosep]
\item $\outp{x}{y} \pref P$ --- send name $y$ over channel $x$.
\item $\inp{x}{y} \pref P$ --- receive a name on channel $x$, bind it to $y$,
  continue as $P$ with $y$ replaced by whatever was received.
\end{itemize}
\end{defn}

\begin{leanbox}
\begin{minted}{lean}
-- π-calculus processes
inductive Pi where
  | nil  : Pi
  | send : Name → Name → Pi → Pi         -- x̄⟨y⟩.P
  | recv : Name → Name → Pi → Pi         -- x(y).P (y is bound)
  | par  : Pi → Pi → Pi                  -- P | Q
  | res  : Name → Pi → Pi                -- (νa)P
  | repl : Pi → Pi                       -- !P
  | plus : Pi → Pi → Pi                  -- P + Q
  deriving Repr

-- Free names
def Pi.freeNames : Pi → List Name
  | .nil => []
  | .send x y P => [x, y] ++ P.freeNames
  | .recv x y P => [x] ++ (P.freeNames.filter (· ≠ y))  -- y is bound
  | .par P Q => P.freeNames ++ Q.freeNames
  | .res a P => P.freeNames.filter (· ≠ a)
  | .repl P => P.freeNames
  | .plus P Q => P.freeNames ++ Q.freeNames
\end{minted}
\end{leanbox}

\section{Communication and Scope Extrusion}

The communication rule: when $\outp{x}{y}.P$ runs in parallel with
$\inp{x}{z}.Q$, they synchronize---$y$ is sent, $z$ gets bound to $y$:
\[
  \outp{x}{y}.P \;\Par\; \inp{x}{z}.Q \;\xrightarrow{\tau}\;
  P \;\Par\; Q[y/z]
\]
Here $Q[y/z]$ is \textbf{substitution}: replace every free occurrence of $z$
in $Q$ by $y$. We give it a Lean definition below (Definition~\ref{def:pi-op}).

The truly novel feature is \textbf{scope extrusion}: if a private name
escapes via communication, its scope automatically widens.
\[
  \res{a}(\outp{x}{a}.P) \;\Par\; \inp{x}{z}.Q
  \;\xrightarrow{\tau}\;
  \res{a}(P \;\Par\; Q[a/z])
  \qquad\text{if } a \notin \fn(Q)
\]
The name $a$ was private to the left process, but after sending it over $x$,
the right process now knows $a$ too. The $\nu a$ scope expands to cover
both processes. The side condition $a \notin \fn(Q)$ is essential: if $a$ were
already free in $Q$ it would be a \emph{different}, unrelated name, and
widening the $\nu a$ scope over $Q$ would wrongly capture it. (This mirrors
the side condition on the structural-congruence version of the same law in
Chapter~7, and in practice one always $\alpha$-renames the bound $a$ to
something fresh for $Q$ first.)

\begin{defn}{$\pi$-Calculus Substitution and Reduction}{pi-op}
Substitution $P[y/z]$ replaces free occurrences of $z$ by $y$, stopping at any
binder that re-binds $z$ (and, by the usual freshness convention, assuming $y$
is not captured by a binder in $P$). The one-step reduction relation
$P \reduces Q$ is generated by communication (in either order) closed under
parallel composition and restriction. The Lean encoding follows.
\end{defn}

\begin{leanbox}
\begin{minted}{lean}
-- Substitution: replace FREE occurrences of z by y  (prose: P[y/z]).
-- Structural recursion on P, so it terminates. Stops at binders that
-- re-bind z; assumes (Barendregt convention) y is fresh for P.
def Pi.subst : Pi → Name → Name → Pi
  | .nil,        _, _ => .nil
  | .send x w P, z, y =>
      .send (if x = z then y else x) (if w = z then y else w) (P.subst z y)
  | .recv x w P, z, y =>
      .recv (if x = z then y else x) w (if w = z then P else P.subst z y)
  | .par P Q,    z, y => .par (P.subst z y) (Q.subst z y)
  | .res a P,    z, y => .res a (if a = z then P else P.subst z y)
  | .repl P,     z, y => .repl (P.subst z y)
  | .plus P Q,   z, y => .plus (P.subst z y) (Q.subst z y)

-- Reduction (τ) semantics. commL/commR: a send meets a matching receive
-- (in either order). parL/parR/res: reduction under those contexts.
inductive PiStep : Pi → Pi → Prop where
  | commL : PiStep (.par (.send x y P) (.recv x z Q)) (.par P (Q.subst z y))
  | commR : PiStep (.par (.recv x z Q) (.send x y P)) (.par (Q.subst z y) P)
  | parL  : PiStep P P' → PiStep (.par P Q) (.par P' Q)
  | parR  : PiStep Q Q' → PiStep (.par P Q) (.par P Q')
  | res   : PiStep P P' → PiStep (.res a P) (.res a P')
\end{minted}
\end{leanbox}

This \lc{PiStep} deliberately keeps only the ``computational'' rules
(communication under parallel/restriction contexts); a fully general
presentation also closes reduction under structural congruence
(Chapter~7), which is what lets scope extrusion fire in arbitrary positions.
We omit that closure to keep the relation small enough to compute with by
hand---it is all we need for the exercises.

\subsection{The Node.js Connection}

This is \emph{exactly} what happens with \lc{MessageChannel} and transferable
objects:

\begin{tsbox}
\begin{minted}{typescript}
// Main thread
const { port1, port2 } = new MessageChannel();  // (νport1)(νport2)

// Send port2 to the worker --- scope extrusion!
worker.postMessage({ port: port2 }, [port2]);

// Now the worker has port2 and can communicate on it directly
// The "private" channel has escaped its original scope
\end{minted}
\end{tsbox}

\section{Encoding the $\lambda$-Calculus}

Milner showed that function application can be encoded as communication.
Each $\lambda$-term is translated relative to a \emph{return channel} $u$ ---
the channel on which its value will be delivered. This is Milner's monadic
call-by-name encoding (every message is a single name, matching our
\lc{send}/\lc{recv} syntax):
\begin{align*}
  ⟦x⟧_u &= \outp{x}{u}.\nil
    &&\text{(look up $x$: hand it the return channel $u$)} \\
  ⟦\lambda x. M⟧_u &=
    \inp{u}{x}.\inp{u}{w}.⟦M⟧_w
    &&\text{(receive the argument channel $x$, then the return channel $w$)} \\
  ⟦M\;N⟧_u &=
    \res{v}\bigl(⟦M⟧_v \;\Par\;
      \outp{v}{n}.\outp{v}{u}.\nil\bigr)
    &&\text{(run $M$ on fresh $v$; send it $N$'s channel, then $u$)}
\end{align*}
where $n$ is the channel giving access to the argument $N$. For a variable
argument $N = y$ this is simply $y$ itself; the general case wraps $N$ in a
replicated server $\res{n}\bigl(\ldots \Par \repl\, \inp{n}{w}.⟦N⟧_w\bigr)$
(see Milner, \emph{Functions as Processes}, 1992).

\begin{warning}[title={Polarity matters}]
The direction of each communication is the whole point. An abstraction
$⟦\lambda x.M⟧_u$ \textbf{inputs} on its channel (it is waiting to be called);
an application \textbf{outputs} on that channel (it is doing the calling). If
both sides input --- as an earlier draft of this encoding had it --- then
nothing ever outputs on $v$, no $\tau$ step is possible, and the ``redex''
$(\lambda x.M)\,N$ is stuck. Getting sender and receiver the right way round
is exactly what makes $\beta$-reduction turn into communication.
\end{warning}

The variable $u$ is a ``return channel''---where the result will be sent.
Function application creates a fresh channel $v$, runs the function on it,
sends the function its argument channel and then the outer return channel,
and the function's body delivers the result on that return channel.

\begin{intuition}
This encoding reveals something deep: sequential computation (the
$\lambda$-calculus) is a special case of concurrent computation (the
$\pi$-calculus) where everything is forced to happen in a specific order
by the channel dependencies. Concurrency is the more fundamental notion;
sequentiality is concurrency with extra constraints.
\end{intuition}

\begin{exercisebox}[title={Exercise 6.1}]
\stepcounter{exercise}
Model a chat room in the $\pi$-calculus: a room process that accepts join
requests, and when two users have joined, creates a private channel between
them. Use scope extrusion to pass the private channel to both users.
\end{exercisebox}

\begin{exercisebox}[title={Exercise 6.2}]
\stepcounter{exercise}
Encode the $\lambda$-term $(\lambda x. x)\; y$ (identity applied to $y$) with
return channel $u$ using the encoding above. You should get
\[
  \res{v}\bigl(\,\inp{v}{x}.\inp{v}{w}.\outp{x}{w}.\nil
    \;\Par\;
    \outp{v}{y}.\outp{v}{u}.\nil\,\bigr).
\]
Now trace the reduction using \lc{PiStep}. Show it takes exactly two
$\tau$ steps --- the first sends $y$ (binding $x := y$), the second sends $u$
(binding $w := u$) --- ending in
$\res{v}\bigl(\outp{y}{u}.\nil \Par \nil\bigr)$, which is structurally
congruent to $⟦y⟧_u = \outp{y}{u}.\nil$ (drop the $\Par\,\nil$ and the vacuous
$\nu v$, since $v \notin \fn$). The result is $y$, as expected.
Encode these processes as \lc{Pi} terms and prove both \lc{PiStep} steps in
Lean by applying \lc{PiStep.res} to \lc{PiStep.commR}. (Worked in the
companion file \lc{PiCalcVerification.lean}.)
\end{exercisebox}

% ══════════════════════════════════════════════════════════════════════════════
\chapter{Structural Congruence and Algebraic Laws}
% ══════════════════════════════════════════════════════════════════════════════

\epigraph{``Algebra is the offer made by the devil to the mathematician.''}%
{\textit{Michael Atiyah}}

\section{The Algebra of Processes}

\begin{historical}[title={Process Algebra}]
Just as high-school algebra studies equations between numerical expressions
($a + b = b + a$, $a \cdot 1 = a$), \textbf{process algebra} studies
equations between concurrent processes. The discovery that processes obey
clean algebraic laws---and that these laws characterize bisimilarity---was
one of the early surprises of concurrency theory. The term ``process algebra''
was coined by Jan Bergstra and Jan Willem Klop in the early 1980s; their
Algebra of Communicating Processes (ACP) developed this algebraic perspective
alongside CCS and CSP.
\end{historical}

\section{Laws of CCS (Up to Bisimilarity)}

\begin{thm}{Algebraic Laws}{ccs-laws}
The following hold for all CCS processes $P$, $Q$, $R$:
\begin{align}
  P \Par \nil &\bisim P
    &&\text{(parallel identity)} \label{eq:par-unit} \\
  P \Par Q &\bisim Q \Par P
    &&\text{(parallel commutativity)} \label{eq:par-comm} \\
  (P \Par Q) \Par R &\bisim P \Par (Q \Par R)
    &&\text{(parallel associativity)} \label{eq:par-assoc} \\
  P + \nil &\bisim P
    &&\text{(choice identity)} \label{eq:plus-unit} \\
  P + Q &\bisim Q + P
    &&\text{(choice commutativity)} \label{eq:plus-comm} \\
  P + P &\bisim P
    &&\text{(choice idempotence)} \label{eq:plus-idem}
\end{align}
\end{thm}

\begin{intuition}
Parallel composition $(\Par, \nil)$ forms a \textbf{commutative monoid}---the
same algebraic structure that addition on integers forms! If you've studied
the Galois theory book, you recognize this: groups, rings, and fields all have
monoid substructure. The algebra of processes and the algebra of numbers share
the same abstract patterns.
\end{intuition}

\section{Structural Congruence}

\begin{defn}{Structural Congruence}{struct-cong}
\textbf{Structural congruence} ($\equiv$) is the smallest equivalence
relation on CCS/\,$\pi$-calculus processes that includes the algebraic laws
above plus:
\begin{align}
  \res{a}\nil &\equiv \nil
    &&\text{(vacuous restriction)} \\
  \res{a}\res{b} P &\equiv \res{b}\res{a} P
    &&\text{(restriction commutes)} \\
  \res{a}(P \Par Q) &\equiv P \Par \res{a}Q
    \quad\text{if } a \notin \fn(P)
    &&\text{(scope extrusion)} \\
  \repl P &\equiv P \Par \repl P
    &&\text{(replication unfolds)}
\end{align}
\end{defn}

\begin{leanbox}
\begin{minted}{lean}
-- Structural congruence as an inductive relation
inductive StructCong : CCS → CCS → Prop where
  | parComm  : StructCong (.par P Q) (.par Q P)
  | parAssoc : StructCong (.par (.par P Q) R) (.par P (.par Q R))
  | parNil   : StructCong (.par P .nil) P
  | plusComm  : StructCong (.plus P Q) (.plus Q P)
  | plusNil   : StructCong (.plus P .nil) P
  | plusIdem  : StructCong (.plus P P) P
  | refl     : StructCong P P
  | symm     : StructCong P Q → StructCong Q P
  | trans    : StructCong P Q → StructCong Q R → StructCong P R
  -- congruence: structural congruence is preserved by all operators
  | parCong  : StructCong P P' → StructCong (.par P Q) (.par P' Q)
  | plusCong : StructCong P P' → StructCong (.plus P Q) (.plus P' Q)
\end{minted}
\end{leanbox}

Note that this Lean \lc{StructCong} encodes only the algebraic laws over the
\lc{CCS} type plus the equivalence-and-congruence closure. It deliberately
omits the four laws in the Definition above that are specific to
restriction and replication --- $\res{a}\nil \equiv \nil$,
$\res{a}\res{b}P \equiv \res{b}\res{a}P$, scope extrusion, and
$\repl P \equiv P \Par \repl P$ --- together with a congruence rule for
$\nu$. Those laws mention constructors (\lc{res}, \lc{repl}) that the
$\pi$-calculus type \lc{Pi} carries; a full \lc{Pi}-level structural
congruence would add them (and a \lc{resCong}). We show the CCS fragment here
because it is all the later CCS material needs; extending it to \lc{Pi} is a
good exercise.

Structural congruence is used to ``prepare'' terms before reduction. The
idea: first rearrange using $\equiv$ (which is purely syntactic), then
apply the communication rule.

\begin{exercisebox}[title={Exercise 7.1}]
\stepcounter{exercise}
Prove in Lean that \lc{StructCong} is an equivalence relation (reflexive,
symmetric, transitive---these are built into the definition above, but
verify that the constructors suffice).
\end{exercisebox}

% ══════════════════════════════════════════════════════════════════════════════
\chapter{Hennessy-Milner Logic}
% ══════════════════════════════════════════════════════════════════════════════

\epigraph{``Logic is the hygiene of the mathematician.''}%
{\textit{André Weil}}

\section{A Logic for Processes}

\begin{historical}[title={Hennessy and Milner (1985)}]
Matthew Hennessy and Robin Milner introduced a modal logic specifically for
reasoning about processes. The logic has two modalities: $\langle a \rangle\varphi$
(``the process \emph{can} do action $a$ and then satisfy $\varphi$'') and
$[a]\varphi$ (``\emph{whenever} the process does $a$, the result satisfies
$\varphi$''). The landmark result: two processes are bisimilar if and only if
they satisfy exactly the same HM formulas. This means bisimulation is not just
one definition of equivalence---it is \emph{the} equivalence that respects all
observable properties.
\end{historical}

\begin{defn}{Hennessy-Milner Logic}{hm-logic}
Given an action set $A$, the formulas of HM logic are:
\[
  \varphi, \psi \;::=\;
    \mathsf{tt}
    \;\mid\; \mathsf{ff}
    \;\mid\; \varphi \wedge \psi
    \;\mid\; \varphi \vee \psi
    \;\mid\; \langle a \rangle \varphi
    \;\mid\; [a] \varphi
\]
\begin{itemize}[nosep]
\item $\langle a \rangle \varphi$ (diamond): ``can do $a$ and then satisfy $\varphi$.''
\item $[a] \varphi$ (box): ``every $a$-step leads to $\varphi$.''
\end{itemize}
\end{defn}

\begin{leanbox}
\begin{minted}{lean}
-- HM Logic formulas
inductive HML (Label : Type) where
  | tt   : HML Label
  | ff   : HML Label
  | conj : HML Label → HML Label → HML Label
  | disj : HML Label → HML Label → HML Label
  | dia  : Action Label → HML Label → HML Label   -- ⟨a⟩φ
  | box  : Action Label → HML Label → HML Label   -- [a]φ
  deriving Repr

-- Satisfaction: when does a state satisfy a formula?
def Satisfies {State Label : Type}
    (step : State → Action Label → State → Prop)
    : State → HML Label → Prop
  | _, .tt => True
  | _, .ff => False
  | s, .conj φ ψ => Satisfies step s φ ∧ Satisfies step s ψ
  | s, .disj φ ψ => Satisfies step s φ ∨ Satisfies step s ψ
  | s, .dia a φ  => ∃ s', step s a s' ∧ Satisfies step s' φ
  | s, .box a φ  => ∀ s', step s a s' → Satisfies step s' φ
\end{minted}
\end{leanbox}

\section{The Adequacy Theorem}

\begin{thm}{HM Logic Characterizes Bisimilarity}{hm-adequacy}
For image-finite LTSs (each state has finitely many transitions):
\[
  s \bisim t \quad\iff\quad
  \forall \varphi \in \text{HML},\;
  s \models \varphi \iff t \models \varphi
\]
\end{thm}

The $\Leftarrow$ direction is proved by showing that ``satisfies the same
formulas'' is itself a bisimulation. The $\Rightarrow$ direction is proved by
induction on the formula structure.

\subsection{Distinguishing the Coffee Machines}

Remember $P$ (choice after coin) and $Q$ (choice before coin) from Chapter~4?
The formula $\langle\text{coin}\rangle(\langle\text{coffee}\rangle\mathsf{tt}
\wedge \langle\text{tea}\rangle\mathsf{tt})$ distinguishes them:
\begin{itemize}[nosep]
\item $P \models$ the formula: after coin, both coffee and tea are possible.
\item $Q \not\models$ the formula: after coin (choosing the coffee branch),
  tea is not possible.
\end{itemize}

\section{Specification via Logic}

HM logic gives us a specification language for concurrent systems:

\begin{itemize}[nosep]
\item ``The server can always accept a request'':
  $[\tau]^*\langle\text{req}\rangle\mathsf{tt}$
  (after any internal steps, a request action is possible).
\item ``After login, logout is always possible'':
  $[\text{login}]\langle\text{logout}\rangle\mathsf{tt}$.
\item ``The system never reaches a deadlocked state'':
  $[a_1]\mathsf{live} \wedge [a_2]\mathsf{live} \wedge \cdots$
  where $\mathsf{live} = \bigvee_a \langle a \rangle \mathsf{tt}$.
\end{itemize}

\begin{warning}[title={Beyond the core logic}]
The first and third examples reach \emph{outside} the HM grammar of
Definition~\ref{def:hm-logic}, and we write them informally. The iterated
modality $[\tau]^*$ (``after any number of $\tau$-steps'') is a Kleene star,
and $\mathsf{live} = \bigvee_a \langle a \rangle\mathsf{tt}$ is a disjunction
over \emph{all} actions --- but the core logic has neither iteration nor
infinite/$n$-ary connectives, only the finite $\wedge$ and $\vee$ of the
grammar. Expressing ``always,'' ``eventually,'' and ``for every reachable
state'' properly requires fixed-point operators: this is the
\emph{modal $\mu$-calculus}, HM logic's more expressive descendant, mentioned
in the next box. Read these two bullets as motivation, not as formulas of the
logic we defined.
\end{warning}

\begin{historical}[title={From HM Logic to Model Checking}]
HM logic is the ancestor of the temporal logics (CTL, CTL*, LTL) used by
industrial model-checking tools like SPIN, NuSMV, and TLA+. Edmund Clarke,
Allen Emerson, and Joseph Sifakis received the 2007 Turing Award for
model checking---the automatic verification of temporal logic properties
against finite-state models. Amazon uses TLA+ to verify protocols in DynamoDB,
S3, and other AWS services.
\end{historical}

\begin{exercisebox}[title={Exercise 8.1}]
\stepcounter{exercise}
Write HM formulas expressing: (a)~``the process can terminate'' and (b)~``the
process is deadlock-free (from every reachable state, some action is
possible).''
\end{exercisebox}

\begin{exercisebox}[title={Exercise 8.2}]
\stepcounter{exercise}
Prove the $\Rightarrow$ direction of the adequacy theorem in Lean: if
$s \bisim t$, then $s \models \varphi \iff t \models \varphi$ for all $\varphi$.
Use induction on $\varphi$.
\end{exercisebox}

% ══════════════════════════════════════════════════════════════════════════════
\part{Types for Concurrency}
% ══════════════════════════════════════════════════════════════════════════════

% ══════════════════════════════════════════════════════════════════════════════
\chapter{Session Types --- Protocols as Types}
% ══════════════════════════════════════════════════════════════════════════════

\epigraph{``Propositions as sessions.''}{\textit{Philip Wadler} (2012)}

\section{Typing Communication}

\begin{historical}[title={Kohei Honda (1956--2012)}]
Kohei Honda had a beautiful idea: if types can describe the structure of
\emph{data} (lists, trees, records), why not the structure of
\emph{communication}? A type like $!\text{Int}.?\text{Bool}.\text{end}$
describes a protocol: ``send an integer, then receive a boolean, then
done.'' The dual type $?\text{Int}.!\text{Bool}.\text{end}$ describes the
other side: ``receive an integer, then send a boolean, then done.''

Honda's 1993 paper introduced \textbf{session types} for the $\pi$-calculus.
The key theorem: if both endpoints of a channel are typed with dual session
types, then communication is \emph{guaranteed} to be deadlock-free,
type-safe, and eventually complete.

Honda (with Vasconcelos and Kubo, 1998) extended the theory to cover choice,
recursion, and delegation (passing a session-typed channel to another
process). Tragically, Honda died in 2012 at the age of 56, before seeing
session types reach mainstream adoption. His work is now recognized as one
of the most important contributions to programming language theory.
\end{historical}

\section{Session Type Syntax}

\begin{defn}{Session Types}{session-types}
\begin{align*}
  S \;::=\;&\; !T.S &&\text{send a value of type $T$, continue as $S$} \\
  \mid\;&\; ?T.S &&\text{receive a value of type $T$, continue as $S$} \\
  \mid\;&\; S_1 \oplus S_2 &&\text{internal choice: sender picks a branch} \\
  \mid\;&\; S_1 \mathbin{\&} S_2 &&\text{external choice: receiver picks a branch} \\
  \mid\;&\; \text{end} &&\text{session complete} \\
  \mid\;&\; \mu X. S &&\text{recursive session (for protocols that loop)}
\end{align*}
\end{defn}

\begin{leanbox}
\begin{minted}{lean}
-- Base data types for messages
inductive BaseType where
  | nat | bool | str
  deriving DecidableEq, Repr

-- Session types
inductive SessionType where
  | send    : BaseType → SessionType → SessionType     -- !T.S
  | recv    : BaseType → SessionType → SessionType     -- ?T.S
  | choose  : SessionType → SessionType → SessionType  -- S₁ ⊕ S₂
  | offer   : SessionType → SessionType → SessionType  -- S₁ & S₂
  | done    : SessionType                              -- end
  | var     : Nat → SessionType                        -- recursion variable
  | mu      : SessionType → SessionType                -- μX.S
  deriving DecidableEq, Repr
\end{minted}
\end{leanbox}

\section{Duality}

The fundamental operation: every session type has a \textbf{dual} that
describes the other endpoint of the channel.

\begin{defn}{Session Type Duality}{duality}
\begin{alignat*}{3}
  \overline{!T.S} &= ?T.\overline{S} &\qquad
  \overline{?T.S} &= !T.\overline{S} \\
  \overline{S_1 \oplus S_2} &= \overline{S_1} \mathbin{\&} \overline{S_2} &\qquad
  \overline{S_1 \mathbin{\&} S_2} &= \overline{S_1} \oplus \overline{S_2} \\
  \overline{\text{end}} &= \text{end} &\qquad
  \overline{\mu X. S} &= \mu X. \overline{S}
\end{alignat*}
\end{defn}

The next listing contains this book's first \emph{tactic} proof: that taking
the dual twice returns the original type. It uses \lc{induction \ldots with}
(one case per constructor, with induction hypotheses \lc{ih} for recursive
sub-types), \lc{simp} (rewrite with the given equations until nothing
changes), and \lc{rfl} (true by computation). If any of these is unfamiliar,
Appendix~\ref{app:primer} explains each one --- including the recursor term
that \lc{induction} is sugar for.

\begin{leanbox}
\begin{minted}{lean}
def SessionType.dual : SessionType → SessionType
  | .send t s    => .recv t s.dual
  | .recv t s    => .send t s.dual
  | .choose a b  => .offer a.dual b.dual
  | .offer a b   => .choose a.dual b.dual
  | .done        => .done
  | .var n       => .var n
  | .mu s        => .mu s.dual

-- Duality is an involution: dual (dual S) = S
theorem SessionType.dual_dual (S : SessionType) :
    S.dual.dual = S := by
  induction S with
  | send t s ih => simp [dual, ih]
  | recv t s ih => simp [dual, ih]
  | choose a b iha ihb => simp [dual, iha, ihb]
  | offer a b iha ihb => simp [dual, iha, ihb]
  | done => rfl
  | var n => rfl
  | mu s ih => simp [dual, ih]
\end{minted}
\end{leanbox}

\section{Type Safety}

The central guarantee: well-typed processes don't deadlock.

\begin{thm}{Session Type Safety}{session-safety}
If process $P$ has session type $S$ on channel $c$, and process $Q$ has
session type $\overline{S}$ on $c$, then $\res{c}(P \Par Q)$ is:
\begin{enumerate}[nosep, label=(\roman*)]
\item \textbf{Type-safe:} messages always have the expected type.
\item \textbf{Deadlock-free:} if both sides follow their types, communication
  always progresses.
\item \textbf{Session-complete:} the protocol eventually reaches $\text{end}$
  on both sides.
\end{enumerate}
\end{thm}

\begin{proof}[Proof sketch]
By \textbf{progress} and \textbf{preservation}:
\begin{itemize}[nosep]
\item \emph{Progress:} if the session hasn't ended, at least one side can
  take a step (because send is matched with receive in the dual).
\item \emph{Preservation:} after a communication step, the remaining session
  types are still dual (because duality ``peels off'' matching prefixes).
\end{itemize}
\end{proof}

\section{The Linear Logic Connection}

\begin{historical}[title={Wadler's ``Propositions as Sessions'' (2012)}]
Philip Wadler discovered a remarkable correspondence: session types are
propositions in \textbf{classical linear logic}, and well-typed processes
are proofs. This is a Curry-Howard correspondence for concurrent
computation:

\begin{center}
\begin{tabular}{lcl}
  \textbf{Linear Logic} & & \textbf{Session Types} \\
  \midrule
  $A \otimes B$ (tensor) & $\longleftrightarrow$ & $!A.S$ (send $A$, continue $S$) \\
  $A \bindnasrepma B$ (par) & $\longleftrightarrow$ & $?A.S$ (receive $A$, continue $S$) \\
  $A \oplus B$ (plus) & $\longleftrightarrow$ & $S_1 \oplus S_2$ (internal choice) \\
  $A \mathbin{\&} B$ (with) & $\longleftrightarrow$ & $S_1 \mathbin{\&} S_2$ (external choice) \\
  $\mathbf{1}$ (unit) & $\longleftrightarrow$ & $\text{end}$ (session complete) \\
  $A^\perp$ (negation) & $\longleftrightarrow$ & $\overline{S}$ (duality) \\
\end{tabular}
\end{center}

The key is \emph{linearity}: each channel is used exactly once (you send on it
or receive on it, then move to the continuation). This is why each message is a
resource consumed exactly once---the linear logic discipline prevents aliasing,
duplication, and forgetting of channels.

If you've read the lambda calculus book: the Curry-Howard correspondence says
``proofs = programs'' for sequential computation. Wadler's correspondence says
``proofs in linear logic = well-typed concurrent programs.''
\end{historical}

\begin{exercisebox}[title={Exercise 9.1}]
\stepcounter{exercise}
Define session types for a two-phase commit protocol: the coordinator sends a
\lc{prepare} message, each participant responds with \lc{vote\_commit} or
\lc{vote\_abort}, then the coordinator sends the final \lc{commit} or
\lc{abort} decision. Verify duality holds.
\end{exercisebox}

\begin{exercisebox}[title={Exercise 9.2}]
\stepcounter{exercise}
Prove \lc{SessionType.dual\_dual} in Lean (the skeleton is above---fill in the
details for each case).
\end{exercisebox}

% ══════════════════════════════════════════════════════════════════════════════
\chapter{Practical Session Types for TypeScript}
% ══════════════════════════════════════════════════════════════════════════════

\epigraph{``In theory, there is no difference between theory and practice. In
practice, there is.''}{\textit{Attributed to Yogi Berra}}

\section{What Session Types Would Look Like}

Imagine if TypeScript could enforce protocol ordering:

\begin{tsbox}
\begin{minted}{typescript}
// A session-typed channel
type WorkerProtocol =
  | { phase: 'send'; payload: Task; next: WorkerRecv }
  | { phase: 'done' };
type WorkerRecv =
  | { phase: 'recv'; payload: Result; next: WorkerDone };
type WorkerDone = { phase: 'done' };

// The type system would enforce:
// 1. You MUST send a Task before receiving a Result
// 2. You MUST receive a Result after sending a Task
// 3. After the exchange, the session is done
\end{minted}
\end{tsbox}

\section{Current TypeScript Approximations}

\subsection{State Machine Pattern}

\begin{tsbox}
\begin{minted}{typescript}
// Encode protocol states as types
type Idle = { state: 'idle' };
type Sent = { state: 'sent'; taskId: string };
type Done = { state: 'done'; result: Result };

// Type-safe transitions
function sendTask(ch: Idle, task: Task): Sent { ... }
function recvResult(ch: Sent): Promise<Done> { ... }
// This won't compile: function recvResult(ch: Idle): ... -- type error!
\end{minted}
\end{tsbox}

This enforces ordering but has limitations: nothing prevents you from using
a channel in state \lc{Sent} twice (non-linearity), or from forgetting to
call \lc{recvResult} (no completion guarantee).

\subsection{Discriminated Unions with XState}

State machine libraries like XState provide runtime enforcement:

\begin{tsbox}
\begin{minted}{typescript}
import { createMachine, interpret } from 'xstate';

const protocolMachine = createMachine({
  initial: 'idle',
  states: {
    idle: { on: { SEND_TASK: 'awaitingResult' } },
    awaitingResult: { on: { RECV_RESULT: 'done' } },
    done: { type: 'final' },
  },
});
\end{minted}
\end{tsbox}

\section{What's Missing in TypeScript}

For full session type safety, TypeScript would need:

\begin{enumerate}[nosep]
\item \textbf{Linear types}: prevent a channel from being used twice or
  forgotten. Rust's ownership system provides this; TypeScript doesn't.
\item \textbf{Higher-kinded types}: abstract over protocol structure
  (e.g., ``for any protocol $S$, ...'').
  This connects back to the lambda calculus book's discussion of System
  F$\omega$.
\item \textbf{Dependent types}: compute the next protocol state from message
  values (e.g., branching on the content of a received message).
\end{enumerate}

\section{Languages with Real Session Type Support}

\begin{itemize}[nosep]
\item \textbf{Rust}: ownership = affine types, so session type libraries
  (\lc{session-types} crate) work well. A consumed channel can't be reused.
\item \textbf{Haskell}: the \lc{simple-sessions} library encodes session
  types using indexed monads.
\item \textbf{Scribble/Links}: research languages designed around multiparty
  session types.
\end{itemize}

\begin{exercisebox}[title={Exercise 10.1}]
\stepcounter{exercise}
Implement a TypeScript wrapper around \lc{Worker} threads that uses the
state machine pattern to enforce a request-response protocol. The API should
make it a type error to receive before sending or to send twice.
\end{exercisebox}

% ══════════════════════════════════════════════════════════════════════════════
\chapter{Multiparty Session Types}
% ══════════════════════════════════════════════════════════════════════════════

\epigraph{``The whole is more than the sum of its parts.''}%
{\textit{Aristotle}}

\section{Beyond Two Parties}

\begin{historical}[title={Honda, Yoshida, and Carbone (2008)}]
Honda, Nobuko Yoshida, and Marco Carbone extended session types from two-party
to \textbf{multiparty} communication. A \emph{global type} describes the
entire protocol from a bird's-eye view; each participant's \emph{local type}
is obtained by \emph{projecting} the global type onto that participant.

The key theorem: if each participant follows their local type, the global
protocol is satisfied. This means you can specify a complex multi-service
interaction once (the global type) and automatically derive each service's
obligations.

The Scribble protocol language (Imperial College London), based on multiparty
session types, has been used by Red Hat, VMware, and the Ocean Observatories
Initiative.
\end{historical}

\section{Global Types}

\begin{defn}{Global Type}{global-type}
A \textbf{global type} specifies the complete interaction:
\[
  G \;::=\;
    p \to q : T \pref G
    \;\mid\; p \to q : \{l_i : G_i\}_{i \in I}
    \;\mid\; \text{end}
    \;\mid\; \mu X. G
\]
\begin{itemize}[nosep]
\item $p \to q : T. G$ --- participant $p$ sends a value of type $T$ to
  participant $q$, then continue as $G$.
\item $p \to q : \{l_i : G_i\}$ --- $p$ sends one of several labels
  $l_i$ to $q$, selecting the corresponding continuation $G_i$.
\end{itemize}
\end{defn}

\subsection{Example: Client-Server-Database}

\begin{align*}
  G = \;&\text{Client} \to \text{Server} : \text{Query}. \\
      &\text{Server} \to \text{Database} : \text{SQL}. \\
      &\text{Database} \to \text{Server} : \text{Rows}. \\
      &\text{Server} \to \text{Client} : \text{Response}. \\
      &\text{end}
\end{align*}

\begin{leanbox}
\begin{minted}{lean}
-- Participants
inductive Role where
  | client | server | database
  deriving DecidableEq, Repr

-- Global types
inductive GlobalType where
  | msg    : Role → Role → BaseType → GlobalType → GlobalType
  | choice : Role → Role → List (String × GlobalType) → GlobalType
  | done   : GlobalType
  | mu     : GlobalType → GlobalType
  | var    : Nat → GlobalType
  deriving Repr
\end{minted}
\end{leanbox}

\section{Projection}

Each participant extracts their local view of the protocol:

\begin{defn}{Projection}{projection}
The \textbf{projection} of global type $G$ onto participant $r$ is:
\begin{align*}
  (p \to q : T. G) \upharpoonright r &=
    \begin{cases}
      !T.(G \upharpoonright r) & \text{if } r = p \\
      ?T.(G \upharpoonright r) & \text{if } r = q \\
      G \upharpoonright r & \text{if } r \neq p, q
    \end{cases}
\end{align*}
\end{defn}

\begin{thm}{Multiparty Session Safety}{multiparty-safety}
If global type $G$ is well-formed (no ambiguity in projections), and each
participant $r$ follows their local type $G \upharpoonright r$, then:
\begin{enumerate}[nosep, label=(\roman*)]
\item Communication is deadlock-free.
\item Messages are type-safe.
\item The protocol completes.
\end{enumerate}
\end{thm}

\section{Connection to Choreographic Programming}

\begin{intuition}
Multiparty session types enable \textbf{choreographic programming}: write the
global protocol once, then \emph{automatically generate} each participant's
code. This is the concurrent analogue of ``write the specification, derive the
implementation''---a theme that connects to formal verification throughout
computer science.

For TypeScript: imagine writing the global protocol for your microservice
architecture and having the TypeScript types for each service's API
automatically derived. gRPC protocol definitions, GraphQL subscriptions, and
WebSocket protocols are all steps toward this vision.
\end{intuition}

\begin{exercisebox}[title={Exercise 11.1}]
\stepcounter{exercise}
Define a global type for an authentication handshake: Client sends credentials
to AuthServer, AuthServer sends a token (or rejection) to Client, on success
Client sends the token to ResourceServer, ResourceServer sends data to Client.
Project onto each of the three roles.
\end{exercisebox}

% ══════════════════════════════════════════════════════════════════════════════
\part{Capstone}
% ══════════════════════════════════════════════════════════════════════════════

% ══════════════════════════════════════════════════════════════════════════════
\chapter{Capstone --- Verified Worker Pool Protocol}
% ══════════════════════════════════════════════════════════════════════════════

\epigraph{``Move fast and break things''\\
\normalfont becomes\\
``Move fast and \emph{prove} things.''}%
{The promise of verified concurrency}

\begin{warning}[title={``Verified'' means two different things --- be precise}]
This is a good moment to be honest about what ``verified'' covers in this book.
There is a firm line between \emph{proved on paper} and \emph{checked by Lean}.
Machine-checked in Lean (the kernel accepts the proof term): the syntax of
every calculus, the session-duality involution
(Theorem/\lc{SessionType.dual\_dual}), the two-step \lc{CCSStep} request/response
derivation, and the two-step \lc{PiStep} reduction of the identity encoding
(Exercise~6.2). Everything with a ``\emph{Proof sketch}'' label below --- session
safety, multiparty safety, and the three worker-pool theorems --- is argued
\emph{on paper} in this chapter, not checked in Lean. Turning those sketches
into machine-checked theorems (which would require, among other things, a
typing relation and a fuel-bounded worker) is exactly the kind of work the
field does and the exercises point toward. Read ``verified'' in the chapter
title as the \emph{goal} the protocol is designed to meet, not a claim that
Lean has certified all of it.
\end{warning}

\section{The Protocol}

We design, formalize, and verify a worker pool protocol for Node.js:

\begin{enumerate}[nosep]
\item A \textbf{Dispatcher} maintains a queue of tasks and a pool of workers.
\item \textbf{Workers} report availability, receive tasks, and return results.
\item Workers can \textbf{join and leave} dynamically ($\pi$-calculus, not
  just CCS).
\end{enumerate}

\subsection{Global Type}

\begin{align*}
  G_{\text{pool}} = \mu X.\; &\text{Worker} \to \text{Dispatcher} :
    \{\text{ready} : G_{\text{task}},\; \text{leave} : \text{end}\} \\
  G_{\text{task}} = \;&\text{Dispatcher} \to \text{Worker} : \text{Task}. \\
    &\text{Worker} \to \text{Dispatcher} : \text{Result}. \\
    &X
\end{align*}

A worker announces it's ready (or leaves). If ready, the dispatcher sends a
task, the worker returns a result, and the protocol loops.

\subsection{$\pi$-Calculus Formalization}

\begin{leanbox}
\begin{minted}{lean}
-- The dispatcher process
def dispatcher (tasks : List Name) (pool : Name) : Pi :=
  -- Listen for worker registrations on pool channel
  .recv pool "worker" <|
  -- Send a task to the worker
  match tasks with
  | [] => .nil  -- no more tasks
  | t :: ts =>
    .send "worker" t <|
    -- Wait for result
    .recv "worker" "result" <|
    dispatcher ts pool

-- A worker process. It loops forever (`worker pool` recurses with the SAME
-- argument and no decreasing measure), so Lean cannot prove it terminates
-- and rejects a plain `def`. That is intentional: a server loop is *meant*
-- to run forever. We mark it `partial def` (which asks Lean to trust us
-- rather than check termination). A `partial def` returning `Pi` also needs
-- to know `Pi` is nonempty, so we register `Inhabited Pi` first.
instance : Inhabited Pi := ⟨.nil⟩

partial def worker (pool : Name) : Pi :=
  -- Register with dispatcher
  .send pool "self" <|       -- scope extrusion: pass our channel
  -- Receive task
  .recv "self" "task" <|
  -- Process and return result (modeled as tau + send)
  Pi.res "result" <|
  .send "self" "result" <|
  worker pool                 -- loop: re-register
\end{minted}
\end{leanbox}

Because \lc{worker} is \lc{partial}, Lean treats it as an opaque black box in
proofs: we cannot unfold it or do induction on its (nonexistent) recursion.
That is a deliberate trade-off --- \lc{partial def} buys us
non-termination at the cost of proof transparency. The \lc{dispatcher} above,
by contrast, recurses structurally on the shrinking \lc{tasks} list, so it is
an ordinary \lc{def} we \emph{can} reason about.

\section{Properties to Prove}

\begin{thm}{Worker Pool Deadlock Freedom}{pool-deadlock-free}
If there is at least one worker and the dispatcher has tasks, the system
always makes progress: there is always a possible transition.
\end{thm}

\begin{proof}[Proof sketch]
\begin{enumerate}[nosep]
\item If a worker is in the ``ready'' state and the dispatcher has a task,
  they can communicate (the ready/task exchange).
\item After receiving a task, the worker eventually produces a result
  (by the structure of the worker process).
\item After receiving the result, the dispatcher loops and the worker
  re-registers.
\item At every state, at least one communication can occur. \qedhere
\end{enumerate}
\end{proof}

\begin{thm}{No Message Loss}{pool-no-loss}
Every task sent by the dispatcher is eventually paired with a result
received from a worker.
\end{thm}

\begin{thm}{Session Type Safety}{pool-type-safe}
Workers always receive well-formed tasks (matching the session type),
and the dispatcher always receives well-formed results.
\end{thm}

\section{The TypeScript Implementation}

\begin{tsbox}
\begin{minted}{typescript}
import { Worker, MessageChannel, MessagePort } from 'worker_threads';

interface Task { id: string; data: unknown; }
interface Result { id: string; output: unknown; }

class VerifiedWorkerPool {
  private workers: Map<string, MessagePort> = new Map();
  private taskQueue: Task[] = [];
  private idle: string[] = [];

  constructor(private workerScript: string, private size: number) {
    for (let i = 0; i < size; i++) this.spawnWorker(`w${i}`);
  }

  private spawnWorker(id: string) {
    const { port1, port2 } = new MessageChannel();  // (ν port1, port2)
    const w = new Worker(this.workerScript);
    w.postMessage({ port: port2 }, [port2]);          // scope extrusion!

    this.workers.set(id, port1);

    port1.on('message', (msg: { type: string } & Record<string, unknown>) => {
      if (msg.type === 'ready') {                     // Worker → Dispatcher: ready
        this.idle.push(id);
        this.dispatch();
      } else if (msg.type === 'result') {             // Worker → Dispatcher: result
        this.handleResult(msg as unknown as Result);
      }
    });
  }

  private dispatch() {
    while (this.idle.length > 0 && this.taskQueue.length > 0) {
      const wid = this.idle.shift()!;
      const task = this.taskQueue.shift()!;
      const port = this.workers.get(wid)!;
      port.postMessage({ type: 'task', ...task });    // Dispatcher → Worker: task
    }
  }

  submit(task: Task) {
    this.taskQueue.push(task);
    this.dispatch();
  }

  private handleResult(result: Result) {
    // Process result...
  }
}
\end{minted}
\end{tsbox}

The comments show exactly where the formal model maps to the implementation:
\lc{MessageChannel()} is restriction ($\nu$), \lc{postMessage} with transfer
is scope extrusion, and the ready/task/result message sequence follows the
session type.

\begin{exercisebox}[title={Exercise 12.1}]
\stepcounter{exercise}
Extend the protocol to handle worker crashes: if a worker fails to return a
result within a timeout, the dispatcher reassigns the task to another worker.
Formalize the extended protocol and prove that every task eventually completes
(assuming at least one worker is always available).
\end{exercisebox}

\begin{exercisebox}[title={Exercise 12.2}]
\stepcounter{exercise}
Argue on paper that the worker pool protocol satisfies deadlock freedom: show
that in every reachable state at least one \lc{PiStep} communication rule
(Definition~\ref{def:pi-op}) is applicable. \emph{Hint:} follow the four-point
proof sketch of Theorem~\ref{thm:pool-deadlock-free}, matching each case to a
\lc{PiStep.commL}/\lc{commR} redex. \emph{Why paper and not Lean here?} Our
\lc{worker} is a \lc{partial def} (it is an infinite loop), so Lean treats it
as opaque: you cannot unfold it or induct on it, which blocks a machine proof.
For a Lean-checkable version, first replace \lc{worker} by a fuel-bounded
structural \lc{def worker (fuel : Nat) (pool : Name) : Pi} that recurses on
\lc{fuel}, then state and prove progress for each fuel step. (Compare the
fully machine-checked two-step \lc{PiStep} reduction of Exercise~6.2.)
\end{exercisebox}

% ══════════════════════════════════════════════════════════════════════════════
\chapter{Looking Back --- The Science of Concurrency}
% ══════════════════════════════════════════════════════════════════════════════

\epigraph{``Computer science is no more about computers than astronomy is
about telescopes.''}{\textit{Edsger W.\ Dijkstra}}

\section{The Intellectual Arc}

We traced the science of concurrency from Dijkstra's identification of the
problem (1965) through Milner's formal languages (CCS 1980, $\pi$-calculus
1992) to Honda's session types (1993) and Wadler's linear logic connection
(2012). Along the way, we formalized every major concept in Lean\,4 and
connected each one to TypeScript/Node.js patterns.

\section{Where the Models Live in Industry}

\begin{center}
\begin{tabular}{lp{8cm}}
  \toprule
  \textbf{Theory} & \textbf{Practice} \\
  \midrule
  CCS / CSP & Go channels, Erlang/Elixir actors, Node.js workers \\
  $\pi$-calculus & Dynamic microservice topologies, Kubernetes service discovery, \lc{MessageChannel()} \\
  Bisimulation & Safe refactoring of concurrent code (replacing one implementation with a bisimilar one) \\
  Session types & gRPC protocol definitions, GraphQL subscriptions, WebSocket protocols \\
  HM logic / temporal logic & TLA+ (Amazon/AWS), SPIN model checker, runtime verification \\
  \bottomrule
\end{tabular}
\end{center}

\section{What We Didn't Cover}

\begin{itemize}[nosep]
\item \textbf{Shared-memory concurrency and separation logic.} Peter O'Hearn
  and John Reynolds developed separation logic for reasoning about programs
  that manipulate shared mutable state (pointers, heaps). O'Hearn received
  the 2016 Gödel Prize for this work. Facebook/Meta uses the Infer tool
  (based on separation logic) to automatically find concurrency bugs in their
  mobile apps.
\item \textbf{Software transactional memory (STM).} Instead of locks, you
  write transactions that are automatically retried on conflict. Haskell's
  STM library is the gold standard; Clojure uses a similar approach.
\item \textbf{The actor model.} Carl Hewitt (1973) proposed actors---objects
  that communicate exclusively via messages. Erlang/OTP is the most
  successful actor-based system. The actor model is closely related to the
  $\pi$-calculus but emphasizes different aspects (named entities vs.\
  anonymous channels).
\item \textbf{Petri nets.} Carl Adam Petri (1962) introduced a graphical
  model of concurrency that predates both CCS and CSP. Petri nets remain
  widely used in manufacturing, logistics, and workflow modeling.
\item \textbf{Async/await.} The \lc{async}/\lc{await} pattern that you use
  daily in TypeScript is syntactic sugar over continuation-passing style
  (CPS), which connects back to the lambda calculus. The delimited
  continuations that underlie \lc{async}/\lc{await} can be given a
  process-calculus interpretation.
\end{itemize}

\section{Cross-Book Connections}

\begin{itemize}[nosep]
\item \textbf{Lambda calculus book:} Sequential computation is a degenerate
  case of concurrent computation. Milner's encoding of the $\lambda$-calculus
  in the $\pi$-calculus makes this precise. Curry-Howard for sequential
  computation (proofs = programs) extends to Wadler's Curry-Howard for
  concurrent computation (proofs in linear logic = session-typed processes).
\item \textbf{Lattice/propagator book:} Propagator networks \emph{are}
  concurrent systems on a lattice. The fixed-point semantics of propagators
  (cells monotonically increase until convergence) is a special case of
  process convergence. Bisimilarity is a greatest fixed point---the same
  Knaster-Tarski theorem.
\item \textbf{Galois theory book:} The algebraic laws of parallel composition
  (commutative monoid) are the same kind of structure as the algebraic laws
  of field arithmetic. The subgroup lattice $\leftrightarrow$ intermediate
  field lattice correspondence is structurally parallel to the
  specification $\leftrightarrow$ implementation correspondence in
  refinement theory.
\end{itemize}

\vfill
\begin{center}
\Large\itshape\color{deepblue}
``The purpose of computing is insight, not numbers.''\\[0.5cm]
\normalsize --- Richard Hamming\\[0.3cm]
\normalsize The purpose of concurrency theory is\\
correctness, not performance.\\
(Performance is a bonus.)
\end{center}
\vfill

% ══════════════════════════════════════════════════════════════════════════════
\backmatter
% ══════════════════════════════════════════════════════════════════════════════

\chapter*{Appendix A: CCS and $\pi$-Calculus Quick Reference}
\addcontentsline{toc}{chapter}{Appendix A: Quick Reference}

\section*{CCS Syntax}

\begin{center}
\begin{tabular}{lcl}
  \toprule
  \textbf{Syntax} & \textbf{Name} & \textbf{Meaning} \\
  \midrule
  $\nil$ & Inaction & Stopped process \\
  $\alpha.P$ & Prefix & Do $\alpha$, continue as $P$ \\
  $P + Q$ & Choice & Behave as $P$ or $Q$ \\
  $P \Par Q$ & Parallel & $P$ and $Q$ run concurrently \\
  $\res{a}P$ & Restriction & Channel $a$ is private \\
  $!P$ & Replication & Unlimited copies of $P$ \\
  \bottomrule
\end{tabular}
\end{center}

\section*{$\pi$-Calculus Extensions}

\begin{center}
\begin{tabular}{lcl}
  \toprule
  \textbf{Syntax} & \textbf{Name} & \textbf{Meaning} \\
  \midrule
  $\outp{x}{y}.P$ & Output & Send $y$ on channel $x$ \\
  $\inp{x}{y}.P$ & Input & Receive on $x$, bind to $y$ \\
  \bottomrule
\end{tabular}
\end{center}

\section*{Session Type Syntax}

\begin{center}
\begin{tabular}{lcl}
  \toprule
  \textbf{Syntax} & \textbf{Name} & \textbf{Dual} \\
  \midrule
  $!T.S$ & Send & $?T.\overline{S}$ \\
  $?T.S$ & Receive & $!T.\overline{S}$ \\
  $S_1 \oplus S_2$ & Internal choice & $\overline{S_1} \mathbin{\&} \overline{S_2}$ \\
  $S_1 \mathbin{\&} S_2$ & External choice & $\overline{S_1} \oplus \overline{S_2}$ \\
  $\text{end}$ & Session complete & $\text{end}$ \\
  \bottomrule
\end{tabular}
\end{center}

\chapter*{Appendix B: Lean\,4 Definitions Index}
\addcontentsline{toc}{chapter}{Appendix B: Lean 4 Definitions Index}

\begin{center}
\begin{tabular}{llr}
  \toprule
  \textbf{Definition} & \textbf{Description} & \textbf{Chapter} \\
  \midrule
  \lc{Action} & Observable or silent action & 2 \\
  \lc{LTS} & Labeled transition system & 2 \\
  \lc{CCS} & CCS process terms & 3 \\
  \lc{CCSAct} & CCS actions (input/output/tau) & 3 \\
  \lc{CCSAct.complement} & Complement of an action & 3 \\
  \lc{CCSStep} & CCS operational semantics & 3 \\
  \lc{IsBisimulation} & Bisimulation condition & 4 \\
  \lc{Bisimilar} & Bisimilarity relation & 4 \\
  \lc{isDeadlocked} & Deadlock predicate & 5 \\
  \lc{philosopher} & Dining-philosopher process & 5 \\
  \lc{Pi} & $\pi$-calculus terms & 6 \\
  \lc{Pi.freeNames} & Free names of a $\pi$-term & 6 \\
  \lc{Pi.subst} & Name substitution $P[y/z]$ & 6 \\
  \lc{PiStep} & $\pi$-calculus reduction semantics & 6 \\
  \lc{StructCong} & Structural congruence & 7 \\
  \lc{HML} & Hennessy-Milner logic formulas & 8 \\
  \lc{Satisfies} & HM logic satisfaction & 8 \\
  \lc{BaseType} & Message base types & 9 \\
  \lc{SessionType} & Session types & 9 \\
  \lc{SessionType.dual} & Session type duality & 9 \\
  \lc{SessionType.dual\_dual} & Duality involution (theorem) & 9 \\
  \lc{Role} & Multiparty participants & 11 \\
  \lc{GlobalType} & Multiparty global types & 11 \\
  \lc{dispatcher} & Worker-pool dispatcher & 12 \\
  \lc{worker} & Worker-pool worker loop & 12 \\
  \bottomrule
\end{tabular}
\end{center}

\chapter*{Appendix C: Historical Timeline}
\addcontentsline{toc}{chapter}{Appendix C: Historical Timeline}

\begin{center}
\begin{tabular}{rl}
  \toprule
  \textbf{Year} & \textbf{Event} \\
  \midrule
  1965 & Dijkstra: ``Cooperating Sequential Processes'' \\
  1971 & Dijkstra: Dining Philosophers \\
  1973 & Hewitt: Actor Model \\
  1976 & Keller: Labeled Transition Systems \\
  1978 & Hoare: Communicating Sequential Processes (paper) \\
  1980 & Milner: CCS \\
  1981 & Park: Bisimulation \\
  1985 & Hennessy \& Milner: HM Logic \\
  1985 & Hoare: CSP (book) \\
  1985--87 & Therac-25 radiation overdoses \\
  1989 & Milner: \textit{Communication and Concurrency} \\
  1992 & Milner, Parrow, Walker: $\pi$-Calculus \\
  1993 & Honda: Session Types \\
  1998 & Honda, Vasconcelos, Kubo: Extended Session Types \\
  2007 & Clarke, Emerson, Sifakis: Turing Award (model checking) \\
  2008 & Honda, Yoshida, Carbone: Multiparty Session Types \\
  2012 & Wadler: ``Propositions as Sessions'' \\
  2026 & You: formalizing it all in Lean\,4 \\
  \bottomrule
\end{tabular}
\end{center}

% ══════════════════════════════════════════════════════════════════════════════
\chapter*{Appendix D: A Lean 4 Primer}\label{app:primer}
\addcontentsline{toc}{chapter}{Appendix D: A Lean 4 Primer}
% ══════════════════════════════════════════════════════════════════════════════

This appendix is a zero-prerequisites tour of every Lean~4 construct and tactic
used in this book. Each entry says what the construct does, shows a small
self-contained example (every example here is machine-checked in the companion
file \lc{PiCalcVerification.lean}), and --- where it illuminates --- shows the
lower-level term the tactic is sugar for. That last part is worth your
attention: Lean tactics are not magic, they are \emph{programs that build proof
terms}, and seeing the term once demystifies the tactic permanently.

Everything here is core Lean~4 (toolchain \lc{leanprover/lean4:v4.31.0}), no
Mathlib --- the same setup as the rest of the book.

\section*{Declarations}

\subsection*{\lc{inductive} --- defining new types}

An \lc{inductive} declaration creates a new type from a list of
\emph{constructors} --- the only ways to build values of the type. The kernel
knows constructors are injective and disjoint, and generates an induction
principle (the \emph{recursor}, below). Every syntax type in this book
(\lc{CCS}, \lc{Pi}, \lc{HML}, \lc{SessionType}, \ldots) is an \lc{inductive}.

\begin{leanbox}
\begin{minted}{lean}
inductive Tree where
  | leaf : Tree
  | node : Tree → Tree → Tree
  deriving Repr, DecidableEq
\end{minted}
\end{leanbox}

\lc{deriving Repr} asks Lean to synthesise a printer (so \lc{\#eval} can show
values); \lc{deriving DecidableEq} synthesises a decidable \lc{=} test (used by
\lc{CCSAct}, \lc{SessionType}, and the \lc{if x = z} tests inside
\lc{Pi.subst}).

\subsection*{\lc{def} and pattern matching}

A \lc{def} introduces a function or value. Functions on inductive types are
usually written by \emph{pattern matching}, one equation per constructor.
Recursion is allowed when some argument structurally shrinks (Lean checks
this).

\begin{leanbox}
\begin{minted}{lean}
def Tree.size : Tree → Nat
  | .leaf     => 1
  | .node l r => l.size + r.size
\end{minted}
\end{leanbox}

The leading dot in \lc{.leaf} means ``the constructor of the type being
matched'' --- shorthand for \lc{Tree.leaf}. Because the function is named
\lc{Tree.size}, \emph{dot notation} \lc{t.size} works for any \lc{t : Tree}.
When Lean \emph{cannot} see that recursion shrinks --- as with the server loop
\lc{worker}, which calls itself with the same argument --- you must write
\lc{partial def} instead (below).

\subsection*{\lc{structure} and \lc{abbrev}}

A \lc{structure} is a one-constructor inductive with named fields (\lc{LTS} in
Chapter~2 is a structure). An \lc{abbrev} is a transparent name for an existing
type (\lc{abbrev Name := String}); unlike \lc{def}, an \lc{abbrev} unfolds
automatically during type-checking.

\begin{leanbox}
\begin{minted}{lean}
structure Point where
  x : Nat
  y : Nat
  deriving Repr

abbrev Nombre := String
\end{minted}
\end{leanbox}

\subsection*{\lc{Prop}, \lc{theorem}, and \lc{example}}

In Lean, propositions are types (of sort \lc{Prop}) and proofs are programs: to
prove $P$ is to construct a term of type $P$. A \lc{theorem} is a named proof;
an \lc{example} is an anonymous one (used throughout this appendix). The body
can be a direct \emph{term} or a \emph{tactic block} introduced by \lc{by}, in
which tactics manipulate the current \emph{goal} until none remain.

\begin{leanbox}
\begin{minted}{lean}
example (p q : Prop) (hp : p) (h : p → q) : q := h hp   -- term style

example (p q : Prop) : p → (p → q) → q := by            -- tactic style
  intro hp h
  apply h
  exact hp
\end{minted}
\end{leanbox}

\subsection*{Type classes and \lc{instance}}

A \emph{type class} is an interface; an \lc{instance} registers an
implementation that Lean finds automatically. The book uses
\lc{instance : Inhabited Pi := ⟨.nil⟩} so that the \lc{partial def worker}
(which returns a \lc{Pi}) is allowed --- a \lc{partial def} must know its
result type is nonempty.

\begin{leanbox}
\begin{minted}{lean}
instance : Inhabited Tree := ⟨.leaf⟩
\end{minted}
\end{leanbox}

The angle brackets \lc{⟨…⟩} are the \emph{anonymous constructor}: Lean infers
which constructor to apply from the expected type. They also build proofs of
$\exists$ and $\land$ --- \lc{⟨k, hk⟩} proves \lc{∃ k, P k} by giving the
witness and its proof, and proves \lc{p ∧ q} by giving both halves.

\subsection*{\lc{partial def} --- loops without a termination proof}

\lc{partial def} tells Lean not to check termination (you assert it, Lean trusts
you). This is how the server loop \lc{worker} is written. The cost: a
\lc{partial def} is \emph{opaque} in proofs --- you cannot unfold it or induct on
it --- which is why the worker-pool deadlock-freedom result stays on paper.

\begin{leanbox}
\begin{minted}{lean}
partial def spin (n : Nat) : Nat := spin n   -- never terminates; legal
\end{minted}
\end{leanbox}

\subsection*{\lc{notation} / \lc{infixl} / \lc{infixr} --- custom syntax}

These declare operators. When marked \lc{scoped}, they \textbf{must} sit inside
a \lc{namespace} (Lean errors otherwise --- exactly the bug fixed in the CCS
notation block of Chapter~3); \lc{open}ing the namespace brings them into scope.
The number is the precedence.

\begin{leanbox}
\begin{minted}{lean}
namespace TreeNotation
scoped notation "◇" => Tree.leaf
scoped infixr:65 " △ " => Tree.node
end TreeNotation
\end{minted}
\end{leanbox}

\subsection*{\lc{@[simp]} --- the simplifier's rule set}

The attribute \lc{@[simp]} marks an equation as a left-to-right rewrite rule
for the \lc{simp} tactic.

\begin{leanbox}
\begin{minted}{lean}
@[simp] theorem size_node (l r : Tree) :
    (Tree.node l r).size = l.size + r.size := rfl

example (t : Tree) : (Tree.node t .leaf).size = t.size + 1 := by
  simp [Tree.size]
\end{minted}
\end{leanbox}

\section*{Core Tactics}

\subsection*{\lc{rfl} --- true by computation}

Proves \lc{a = a} \emph{after both sides compute} to the same normal form.
Definitional unfolding is free: \lc{2 + 3 = 5} is \lc{rfl}. This is what closes
the base cases of \lc{dual\_dual}.

\begin{leanbox}
\begin{minted}{lean}
example : 2 + 3 = 5 := rfl
\end{minted}
\end{leanbox}

\subsection*{\lc{intro}, \lc{exact}, \lc{apply}}

If the goal is \lc{p → q} (or \lc{∀ x, P x}), \lc{intro hp} moves the antecedent
into the context as \lc{hp} --- the tactic form of writing \lc{fun hp => …}.
\lc{exact h} closes the goal with the term \lc{h}. \lc{apply h} works
\emph{backwards}: if \lc{h : p → q} and the goal is \lc{q}, then \lc{apply h}
leaves the new goal \lc{p}. The tactic proof of \lc{p → (p → q) → q} above
assembles exactly the term-style proof \lc{h hp}.

\subsection*{\lc{refine} --- a term with named holes}

\lc{refine e} is like \lc{exact e} but lets you leave holes \lc{?\_} to be
discharged as further goals --- ideal for building an \lc{⟨…, …⟩} proof one part
at a time.

\begin{leanbox}
\begin{minted}{lean}
example (p q : Prop) (hp : p) (hq : q) : p ∧ q := by
  refine ⟨?_, ?_⟩
  · exact hp
  · exact hq
\end{minted}
\end{leanbox}

\subsection*{\lc{subst} --- substitute an equation}

Given \lc{h : a = b}, \lc{subst h} replaces one variable by the other
everywhere and discards \lc{h}.

\begin{leanbox}
\begin{minted}{lean}
example (a b : Nat) (h : a = b) : b = a := by
  subst h
  rfl
\end{minted}
\end{leanbox}

\subsection*{\lc{simp} --- rewrite to a normal form}

\lc{simp} rewrites with its rule set (and any lemmas you pass in brackets)
until nothing changes. \lc{simp [dual, ih]} in \lc{dual\_dual} unfolds \lc{dual}
and applies the induction hypothesis. When a goal reduces to arithmetic, the
core tactics \lc{omega} (linear arithmetic over integers/naturals) and
\lc{decide} (decidable propositions) finish it.

\section*{Induction and Case Analysis}\label{sec:primer-induction}

\subsection*{\lc{cases} --- split on a constructor}

\lc{cases t with} produces one goal per constructor of \lc{t}'s type, with no
induction hypotheses.

\begin{leanbox}
\begin{minted}{lean}
example (t : Tree) : t.size = t.size := by
  cases t with
  | leaf => rfl
  | node l r => rfl
\end{minted}
\end{leanbox}

\subsection*{\lc{induction … with} --- and the recursor it desugars to}

\lc{induction t with} is like \lc{cases}, but each recursive constructor also
gives you an \emph{induction hypothesis} (named \lc{ih} here) --- the statement
already proved for the sub-terms. This is the tactic used in \lc{dual\_dual}.

\begin{leanbox}
\begin{minted}{lean}
theorem size_pos (t : Tree) : 0 < t.size := by
  induction t with
  | leaf => decide
  | node l r ihl ihr => simp [Tree.size]; omega
\end{minted}
\end{leanbox}

\noindent\textbf{Under the hood.} \lc{induction} is sugar for applying the
type's auto-generated \emph{recursor} \lc{Tree.rec}. The recursor takes one
argument per constructor (the recursive ones receive the results for the
sub-terms) plus a \lc{motive} --- the property being proved, as a function of
the value. The identical proof written as a raw term:

\begin{leanbox}
\begin{minted}{lean}
theorem size_pos' (t : Tree) : 0 < t.size :=
  Tree.rec (motive := fun t => 0 < t.size)
    (by decide)                                  -- leaf case
    (fun l r ihl ihr => by simp [Tree.size]; omega) -- node case
    t
\end{minted}
\end{leanbox}

Seeing this once explains why \lc{induction} gives you exactly those cases and
exactly those hypotheses: they are the arguments of \lc{Tree.rec}.

\section*{Structured Terms and Sugar}

\subsection*{\lc{obtain ⟨…⟩} --- destructure a hypothesis}

\lc{obtain ⟨n, hn⟩ := h} unpacks an \lc{∃} or \lc{∧} hypothesis into its
components. It is sugar for a \lc{match} (equivalently \lc{Exists.elim}): both
term forms below type-check.

\begin{leanbox}
\begin{minted}{lean}
example (h : ∃ n : Nat, n = 3) : True := by
  obtain ⟨n, hn⟩ := h
  trivial

example (h : ∃ n : Nat, n = 3) : True :=
  match h with
  | ⟨_, _⟩ => trivial
\end{minted}
\end{leanbox}

\subsection*{\lc{calc} --- chained (in)equalities}

\lc{calc} writes a chain of steps, each justified by a proof; Lean glues them
with \lc{Trans.trans} (transitivity). This is how the multi-step reductions of
the book read on paper.

\begin{leanbox}
\begin{minted}{lean}
example (a b c : Nat) (h1 : a = b) (h2 : b = c) : a = c :=
  calc a = b := h1
    _ = c := h2
\end{minted}
\end{leanbox}

\subsection*{Proof by contradiction --- \lc{Classical.byContradiction}}

To prove \lc{p}, assume \lc{¬p} and derive \lc{False}. Core Lean provides this
as the term \lc{Classical.byContradiction}; Mathlib wraps it in a \lc{by\_contra}
tactic, but this book stays core-only, so we use the term.

\begin{leanbox}
\begin{minted}{lean}
example (p : Prop) (h : ¬ ¬ p) : p :=
  Classical.byContradiction (fun hnp => h hnp)
\end{minted}
\end{leanbox}

\subsection*{\lc{do}-notation --- sugar for monadic bind}

The \lc{do} block sequences monadic computations; each \lc{let x ← e}
desugars to a \lc{bind} (\lc{>>=}). The two definitions below are equal by
\lc{rfl}. (The book's \lc{s!"pick\{i\}"} string interpolation is a similar piece
of sugar, expanding to \lc{toString} calls and concatenation.)

\begin{leanbox}
\begin{minted}{lean}
def firstTwo (xs : List Nat) : Option (Nat × Nat) := do
  let a ← xs[0]?
  let b ← xs[1]?
  pure (a, b)

def firstTwo' (xs : List Nat) : Option (Nat × Nat) :=
  xs[0]? >>= fun a => xs[1]? >>= fun b => pure (a, b)
\end{minted}
\end{leanbox}

\subsection*{\lc{sorry} --- an admitted hole}

\lc{sorry} closes any goal without a proof. It type-checks but is flagged with a
warning; the book uses it for exactly one exercise (the deadlock of
Chapter~5). Never ship a \lc{sorry} as a real proof.

\end{document}
